% 历史模板备份，不是项目编译入口；正式入口为 main.tex。
\documentclass[11pt,a4paper,oneside,openany]{ctexbook}

% -------------------------------------------------
% 页面与正文
% -------------------------------------------------
\usepackage[
  top=2.3cm,
  bottom=2.3cm,
  left=2.45cm,
  right=2.45cm,
  headheight=15pt
]{geometry}

\usepackage{setspace}
\setstretch{1.24}

\setlength{\parindent}{2em}
\setlength{\parskip}{0.25em}

% -------------------------------------------------
% 颜色
% -------------------------------------------------
\usepackage{xcolor}

\definecolor{SoftBlue}{RGB}{76,112,150}
\definecolor{SoftGray}{RGB}{245,247,250}
\definecolor{BlockBlue}{RGB}{220,232,245}
\definecolor{BlockBlueDark}{RGB}{200,220,240}
\definecolor{BaseGray}{RGB}{135,145,158}
\definecolor{COPFBlue}{RGB}{60,102,170}
\definecolor{WarmOrange}{RGB}{210,126,71}
\definecolor{SoftGreen}{RGB}{83,145,126}
\definecolor{DarkText}{RGB}{28,31,35}
\definecolor{LightTan}{RGB}{247,243,236}

\color{DarkText}

% -------------------------------------------------
% 数学、表格、图片
% -------------------------------------------------
\usepackage{amsmath,amssymb,mathtools}
\usepackage{graphicx}
\graphicspath{{figures/}}

\usepackage{booktabs}
\usepackage{tabularx}
\usepackage{longtable}
\usepackage{array}
\usepackage{multirow}

\newcolumntype{Y}{>{\raggedright\arraybackslash}X}
\newcolumntype{C}{>{\centering\arraybackslash}X}

% -------------------------------------------------
% 图形
% -------------------------------------------------
\usepackage{tikz}
\usetikzlibrary{
  arrows.meta,
  positioning,
  calc,
  fit,
  backgrounds,
  shapes.geometric
}

\usepackage{pgfplots}
\pgfplotsset{compat=1.18}

\tikzset{
  myarrow/.style={
    -{Latex[length=2.4mm]},
    line width=1.15pt,
    draw=SoftBlue
  },
  softarrow/.style={
    -{Latex[length=2.1mm]},
    line width=0.9pt,
    draw=gray!55
  },
  card/.style={
    rounded corners=5pt,
    draw=SoftBlue!38,
    fill=SoftGray,
    line width=0.8pt,
    inner sep=6pt
  },
  bluecard/.style={
    rounded corners=5pt,
    draw=COPFBlue!55,
    fill=BlockBlue,
    line width=0.95pt,
    inner sep=6pt
  },
  orangecard/.style={
    rounded corners=5pt,
    draw=WarmOrange!70!black,
    fill=WarmOrange!8,
    line width=0.95pt,
    inner sep=6pt
  },
  greencard/.style={
    rounded corners=5pt,
    draw=SoftGreen!70!black,
    fill=SoftGreen!8,
    line width=0.95pt,
    inner sep=6pt
  }
}

% -------------------------------------------------
% 提示框
% -------------------------------------------------
\usepackage[most]{tcolorbox}

\tcbset{
  lecturebox/.style={
    enhanced,
    breakable,
    boxrule=0.8pt,
    arc=3pt,
    left=7pt,
    right=7pt,
    top=6pt,
    bottom=6pt,
    before skip=10pt,
    after skip=10pt,
    fonttitle=\bfseries
  }
}

\newtcolorbox{principlebox}[1][]{
  lecturebox,
  colback=BlockBlue!68,
  colframe=SoftBlue,
  coltitle=DarkText,
  title={通用原则},
  #1
}

\newtcolorbox{teachingbox}[1][]{
  lecturebox,
  colback=SoftGray,
  colframe=BaseGray!55,
  coltitle=DarkText,
  title={教学简化},
  #1
}

\newtcolorbox{warningbox}[1][]{
  lecturebox,
  colback=WarmOrange!8,
  colframe=WarmOrange,
  coltitle=DarkText,
  title={常见错误},
  #1
}

\newtcolorbox{confirmbox}[1][]{
  lecturebox,
  colback=WarmOrange!6,
  colframe=WarmOrange!85!black,
  coltitle=DarkText,
  title={需实际确认},
  #1
}

\newtcolorbox{practicebox}[1][]{
  lecturebox,
  colback=SoftGreen!8,
  colframe=SoftGreen,
  coltitle=DarkText,
  title={正确实践},
  #1
}

\newtcolorbox{casebox}[1][]{
  lecturebox,
  colback=LightTan,
  colframe=WarmOrange!55,
  coltitle=DarkText,
  title={贯穿案例},
  #1
}

\newtcolorbox{outputbox}[1][]{
  lecturebox,
  colback=SoftGreen!7,
  colframe=SoftGreen!85!black,
  coltitle=DarkText,
  title={本章产出},
  #1
}

\newtcolorbox{quizbox}[1][]{
  lecturebox,
  colback=SoftGray,
  colframe=SoftBlue!65,
  coltitle=DarkText,
  title={章节自测},
  #1
}

% -------------------------------------------------
% 代码
% -------------------------------------------------
\usepackage{listings}

\lstdefinestyle{pythonstyle}{
  language=Python,
  basicstyle=\ttfamily\small,
  backgroundcolor=\color{SoftGray},
  keywordstyle=\color{COPFBlue}\bfseries,
  commentstyle=\color{SoftGreen!70!black},
  stringstyle=\color{WarmOrange!85!black},
  numbers=left,
  numberstyle=\tiny\color{BaseGray},
  numbersep=8pt,
  frame=single,
  rulecolor=\color{BaseGray!35},
  breaklines=true,
  showstringspaces=false,
  tabsize=4,
  columns=fullflexible,
  captionpos=b
}

\lstset{style=pythonstyle}

% -------------------------------------------------
% 列表与图注
% -------------------------------------------------
\usepackage{enumitem}

\setlist[itemize]{
  leftmargin=2em,
  itemsep=0.25em,
  topsep=0.35em
}

\setlist[enumerate]{
  leftmargin=2.2em,
  itemsep=0.3em,
  topsep=0.35em
}

\usepackage{caption}
\captionsetup{
  font=small,
  labelfont=bf,
  labelsep=quad
}

% -------------------------------------------------
% 章节标题
% -------------------------------------------------
\ctexset{
  chapter={
    format=\raggedright\Huge\bfseries\color{DarkText},
    nameformat=\color{SoftBlue},
    numberformat=\color{SoftBlue},
    titleformat=\color{DarkText},
    beforeskip=-5pt,
    afterskip=24pt
  },
  section={
    format=\Large\bfseries\color{SoftBlue},
    beforeskip=1.4ex,
    afterskip=0.8ex
  },
  subsection={
    format=\large\bfseries\color{DarkText},
    beforeskip=1.2ex,
    afterskip=0.55ex
  }
}

\newcommand{\keyterm}[1]{\textbf{\color{SoftBlue}#1}}
\newcommand{\riskterm}[1]{\textbf{\color{WarmOrange!90!black}#1}}
\newcommand{\goodterm}[1]{\textbf{\color{SoftGreen!80!black}#1}}

% -------------------------------------------------
% 页眉页脚
% -------------------------------------------------
\usepackage{fancyhdr}

\pagestyle{fancy}
\fancyhf{}
\fancyhead[L]{\small\color{BaseGray}\leftmark}
\fancyhead[R]{\small\color{BaseGray}量化研究入门}
\fancyfoot[C]{\small\color{BaseGray}\thepage}
\renewcommand{\headrulewidth}{0.4pt}
\renewcommand{\headrule}{
  \hbox to\headwidth{
    \color{SoftBlue!45}\leaders\hrule height \headrulewidth\hfill
  }
}

% -------------------------------------------------
% 链接与引用
% -------------------------------------------------
\usepackage[
  colorlinks=true,
  linkcolor=SoftBlue,
  urlcolor=SoftBlue,
  citecolor=SoftGreen!70!black
]{hyperref}

\usepackage[nameinlink,noabbrev]{cleveref}

% -------------------------------------------------
% 文档信息
% -------------------------------------------------
\title{
  \textbf{A 股日频横截面多因子研究入门}\\[0.6em]
  \large 从点时数据到含成本组合回测
}
\author{量化策略团队新人培训资料}
\date{\today}

\begin{document}

\frontmatter

\maketitle
\tableofcontents

\chapter*{阅读说明}
\addcontentsline{toc}{chapter}{阅读说明}

本讲义面向具备基础金融知识和 Python 数据分析能力，
但尚未独立完成过多因子研究的新人。

\mainmatter

\chapter{研究问题与可证伪假设}
\label{ch:hypothesis}

\section{本章导读}

量化研究的起点不是因子公式，也不是一条回测曲线，而是一个允许失败的研究问题。新人常从“价格趋势可能延续”或“盈利能力较高的公司可能更受市场认可”出发；这些说法提供了直觉，却没有说明用什么历史信息、在什么股票池、预测哪个未来区间，也没有规定什么证据会让研究者放弃原想法。若这些边界在看到结果以后才补写，研究就很容易退化为寻找漂亮曲线。

本章把模糊想法逐步转换为可登记、可复现、可支持、可否定或可判定为证据不足的任务。贯穿案例仍是“宽基股票池三因子周频约束多头策略”：候选因子为中期动量、账面市值比价值和 ROE 类质量，最终可能形成不加杠杆、不做空的约束型多头组合，但本章不计算因子表现，也不预设任何因子有效。

\begin{principlebox}
\textbf{[通用原则]} 一个合格假设必须允许不利结论。正收益不是课程验收条件；发现时间边界错误、证据不足或假设不受支持，同样是有价值且必须保留的研究产出。
\end{principlebox}

\section{学习目标}

完成本章后，读者应能够：

\begin{enumerate}
  \item 区分投资直觉、经济机制、可计算代理与统计假设；
  \item 写清预测对象、方向、股票池、信号时点、持有期和评价口径；
  \item 写出原假设、备择假设以及能够否定想法的证据；
  \item 区分主分析、预注册敏感性分析与探索性分析；
  \item 在实验前冻结主参数、停止条件与版本身份；
  \item 为三个教学因子建立研究假设卡；
  \item 接受支持、带条件支持、不支持和证据不足都是合法结论；
  \item 使用研究编号与实验编号保留完整审计链。
\end{enumerate}

\section{量化研究不是寻找漂亮曲线}

研究者可以轻易改变回看窗口、财务滞后、股票池、持有期或样本起点。若每次变化后都查看结果，再只保留最有利版本，那么最后一条曲线同时包含了投资想法和研究者的选择，无法再被当作独立证据。更危险的是，结果往往会反向改写故事：先找到相关性，再补写一个似乎合理的机制。

正确目标是建立一条可核对的证据链：先说明为什么可能存在预测关系，再冻结代理、时间边界与评价规则，最后让数据决定结论。经济逻辑不是因子已经有效的证明；统计结果也不能创造原本不存在的经济逻辑。

\begin{warningbox}
先批量试参数、只展示最佳结果，再补写研究动机，会隐藏研究自由度并夸大证据。即使代码和每一次计算都正确，选择过程仍会损害结论的可信度。
\end{warningbox}

\section{从投资想法到证据链}

\cref{fig:ch01-research-chain} 给出全书的研究地图。每一步只承担一种职责：经济机制解释“为什么可能成立”，点时数据回答“当时能否知道”，单因子与稳健性证据回答“是否值得保留”，组合和回测才回答“能否转化为约束下的实际持仓与含成本结果”。后一步不能回头改写前一步的冻结定义。

\begin{figure}[htbp]
  \centering
  \begin{tikzpicture}[
    node distance=6mm and 7mm,
    every node/.style={font=\footnotesize, align=center},
    quantHypothesisFlow/.style={quantCard, text width=3.2cm, minimum height=0.95cm},
    quantHypothesisGate/.style={quantBlueCard, text width=3.2cm, minimum height=0.95cm},
    quantHypothesisOutput/.style={quantGreenCard, text width=3.2cm, minimum height=0.95cm}
  ]
    \node[quantHypothesisFlow] (intuition) {投资直觉};
    \node[quantHypothesisFlow, right=of intuition] (mechanism) {经济机制};
    \node[quantHypothesisFlow, right=of mechanism] (proxy) {可计算代理};
    \node[quantHypothesisGate, below=of proxy] (pit) {点时数据};
    \node[quantHypothesisGate, left=of pit] (single) {单因子证据};
    \node[quantHypothesisGate, left=of single] (robust) {稳健性};
    \node[quantHypothesisOutput, below=of robust] (multi) {多因子};
    \node[quantHypothesisOutput, right=of multi] (portfolio) {约束组合};
    \node[quantHypothesisOutput, right=of portfolio] (backtest) {含成本回测};
    \draw[quantArrow] (intuition) -- (mechanism);
    \draw[quantArrow] (mechanism) -- (proxy);
    \draw[quantArrow] (proxy) -- (pit);
    \draw[quantArrow] (pit) -- (single);
    \draw[quantArrow] (single) -- (robust);
    \draw[quantArrow] (robust) -- (multi);
    \draw[quantArrow] (multi) -- (portfolio);
    \draw[quantArrow] (portfolio) -- (backtest);
  \end{tikzpicture}
  \caption{从投资直觉到含成本回测的全文研究流程。箭头表示证据与版本的前向依赖，不表示任何教学因子已经通过。}
  \label{fig:ch01-research-chain}
\end{figure}

其中，可计算代理必须连接到明确的预测标签；检验设计至少覆盖有效覆盖率、RankIC 方向、分组单调性、非目标暴露、换手、成本与稳健性。只有完成这些检查后，候选才可能进入第 \cref{ch:factor-combination,ch:portfolio,ch:backtest} 的多因子、组合和回测阶段。

\section{明确研究任务}

投资想法、研究假设、因子公式和交易规则容易被混用，但它们回答的问题不同。\cref{tab:ch01-object-differences} 是第一次登记时的最小辨别表。

\begin{table}[htbp]
  \centering
  \small
  \caption{投资想法、研究假设、因子公式与交易规则的区别}
  \label{tab:ch01-object-differences}
  \begin{tabularx}{\textwidth}{p{2.4cm} p{4.2cm} Y Y}
    \toprule
    对象 & 回答的问题 & 合格示例的结构 & 不能替代 \\
    \midrule
    投资想法 & 哪种现象值得研究 & “价格趋势可能延续” & 机制、数据边界和检验 \\
    研究假设 & 在什么冻结条件下，什么证据支持或反对预测关系 & 股票池、时点、标签、方向、证据维度和停止条件 & 已实现公式或交易结论 \\
    因子公式 & 如何把历史可得信息转为逐证券数值 & 输入字段、窗口、缺失、方向和版本 & 假设成立或可交易性 \\
    交易规则 & 如何把冻结信号转为订单与持仓 & 选股、权重、约束、成交和成本规则 & 因子预测力与稳健性 \\
    \bottomrule
  \end{tabularx}
\end{table}

研究问题可以写成一句完整问句：“在冻结的历史宽基股票池与周度标签下，某一候选代理是否提供方向一致、可解释、具有增量且具备基本可交易性的横截面预测信息？”其中“某一候选代理”必须替换为版本化定义；“冻结”表示不能看完结果后无痕改动。

\section{预测对象与未来收益标签}

预测对象不是模糊的“未来表现”，而是信号时点之后、由两个可配置成交代理时点界定的横截面未来收益。设证券 \(i\) 在信号日 \(t\) 的因子值为 \(x_{i,t}\)，下一次和再下一次假设成交价格分别为 \(P^{\mathrm{exec}}_{i,t+1}\) 与 \(P^{\mathrm{exec}}_{i,t+1+H}\)，教学标签写为

\begin{equation}
  y_{i,t}^{(H)}
  = \frac{P^{\mathrm{exec}}_{i,t+1+H}}
           {P^{\mathrm{exec}}_{i,t+1}} - 1.
  \label{eq:ch01-forward-label}
\end{equation}

公式只规定时间方向，不替真实项目决定开盘价、收盘价、均价、复权和不可成交处理。第 \cref{ch:universe-labels} 将正式定义标签起止与失败状态；本章只冻结“信号截止严格早于标签起点”，并禁止用标签期信息选择因子或样本。

\begin{confirmbox}
\textbf{[需实际确认]} 成交价格代理、持有期交易日计数、复权与公司行动、缺失成交、停牌和价格限制的处理，均依赖真实数据与研究制度。它们必须先登记，再由后续章节实现。
\end{confirmbox}

\section{股票池、信号时点与持有期}

贯穿案例使用有历史成分记录、可通过配置替换的宽基股票池。每周最后一个交易日收盘后冻结信号，下一交易日按配置代理模拟成交，持有至下一次调仓对应的成交时点。信号日股票池只能由当时及以前的信息确定；下一交易日出现的停牌或价格限制属于成交事件，不能回头删除信号日证券。

股票池、信号频率、调仓频率和持有期共同决定“比较谁、在何时预测、预测多远”。第 \cref{ch:data-timing} 处理数据可得性，第 \cref{ch:universe-labels} 处理历史成分、可交易边界和标签；本章登记这些接口，不提前重复实现。

\section{经济机制与可计算代理}

一条研究链至少包含四个相互独立的陈述：投资直觉提出可能现象；经济机制说明为何历史信息可能缓慢进入价格；可计算代理把机制映射为当时可得的数值；统计假设规定在冻结评价下什么证据支持或反对预测关系。

例如，“趋势可能延续”是直觉；信息扩散缓慢、投资者反应不足或机构交易延续是候选机制；跳过近期区间的中期历史收益是代理。机制可能不完整，代理也可能带有波动或流动性暴露。因此研究者必须同时记录代理为何与机制相关，以及它还可能测量了什么。第 \cref{ch:factor-construction,ch:preprocessing} 分别核对公式实现和基础暴露控制。

\begin{teachingbox}
\textbf{[教学简化]} 机制的作用是形成可检查预测并解释失败路径，不是要求新人证明唯一因果关系。一个逻辑合理但数据不支持的假设应如实停止；一个数据相关但缺乏可解释机制的发现只能进入新一轮探索。
\end{teachingbox}

\section{原假设、备择假设与可证伪性}

在因子方向统一为“数值越高，预期未来相对收益越高”后，一般形式可以写为 \cref{tab:ch01-hypotheses}。它们针对一组预注册证据，而不是单个平均指标。

\begin{table}[htbp]
  \centering
  \small
  \caption{原假设、备择假设与所需证据}
  \label{tab:ch01-hypotheses}
  \begin{tabularx}{\textwidth}{p{1.5cm} p{7.1cm} Y}
    \toprule
    假设 & 一般表述 & 评价纪律 \\
    \midrule
    \(H_0\) & 在冻结股票池、标签和评价口径下，因子不提供稳定、具有增量且具备基本可交易性的横截面预测信息。 & 允许由无效、方向不稳、非目标暴露、换手成本压力或缺乏增量等多种证据支持。 \\
    \(H_1\) & 在相同冻结条件下，因子在预注册证据维度中提供方向一致、可解释并具有基本可交易性的横截面预测信息。 & 需要覆盖率、方向、单调性、暴露、换手、成本、稳健性与增量证据共同支撑。 \\
    \bottomrule
  \end{tabularx}
\end{table}

“平均 IC 大于零”不能单独定义 \(H_1\)：平均值可能被少数时期驱动，也未说明覆盖、分组结构、风险暴露和交易压力。未拒绝 \(H_0\) 不等于证明因子永远无效；证据不足也不等于通过或失败。具体统计阈值、最小有效期数与业务闸门均须在真实项目中事前确认，本讲义不虚构统一标准。

\section{主分析、敏感性分析与探索性分析}

三类分析的区别不在于代码函数，而在于何时定义、承担何种结论。预注册敏感性只检查少量合理邻域；看到结果后产生的新设定属于探索性分析，必须进入新版本。

\begin{table}[htbp]
  \centering
  \footnotesize
  \setlength{\tabcolsep}{3pt}
  \caption{主分析、预注册敏感性分析与探索性分析}
  \label{tab:ch01-analysis-types}
  \begin{tabularx}{\textwidth}{p{2.0cm} p{2.1cm} p{2.5cm} p{2.5cm} p{1.6cm} Y}
    \toprule
    类型 & 定义时间 & 参数范围 & 能否影响本轮主结论 & 是否新版本 & 报告方式 \\
    \midrule
    主分析 & 查看本轮结果前 & 唯一冻结基线 & 是，承担主要结论 & 初始冻结版本 & 完整报告全部证据维度 \\
    预注册敏感性 & 查看本轮结果前 & 少量、有经济或实现理由的邻近设定 & 只能评价依赖性与适用边界，不能替换主分析选赢家 & 每个候选独立配置身份 & 与主分析并列报告，包括失败 \\
    探索性分析 & 查看结果后产生 & 新代理、新窗口或新分组 & 否，只能生成下一轮假设 & 必须新建研究或实验版本 & 标记探索性，保留产生背景 \\
    \bottomrule
  \end{tabularx}
\end{table}

\section{参数预注册与研究自由度}

预注册是在计算结果前记录主窗口、替代窗口、数据滞后、标签、预处理、评价指标、样本切分和停止条件。它不要求研究者预知所有问题，而是让后来者看见哪些选择是事前决定，哪些是看到结果后的调整。

研究自由度包括公式、窗口、缺失处理、异常值处理、行业控制、股票池、标签、样本起止和汇报指标。候选越多，偶然出现有利结果的机会越大。主配置应保持唯一；敏感性集合应小且有理由；超出集合的新想法要创建新版本，不得覆盖旧记录。第 \cref{ch:robustness} 将进一步管理时间切分、参数邻域和测试集访问。

\section{支持、否定、返工和证据不足}

证据结论与工作流状态应分开表达。证据可以“支持”“带条件支持”“不支持”或“证据不足”；工作流则可以继续、带条件继续、返回实现或数据修复、停止或保留证据不足。返工只适用于数据或实现错误，不用于挽救可信但不利的结果。

\begin{figure}[htbp]
  \centering
  \begin{tikzpicture}[
    node distance=6mm and 10mm,
    every node/.style={font=\footnotesize, align=center},
    quantHypothesisDecision/.style={diamond, aspect=2.0, draw=SoftBlue, fill=BlockBlue, inner sep=1.2pt, text width=2.9cm},
    quantHypothesisContinue/.style={quantGreenCard, text width=3.0cm},
    quantHypothesisRepair/.style={quantOrangeCard, text width=3.0cm},
    quantHypothesisInsufficient/.style={quantCard, text width=3.0cm}
  ]
    \node[quantHypothesisDecision] (valid) {数据与实现可信？};
    \node[quantHypothesisRepair, left=of valid] (repair) {返回实现或数据修复\\新版本保留旧记录};
    \node[quantHypothesisDecision, below=of valid] (enough) {证据足以评价？};
    \node[quantHypothesisInsufficient, left=of enough] (insufficient) {证据不足\\不通过也不失败};
    \node[quantHypothesisDecision, below=of enough] (support) {预注册证据支持？};
    \node[quantHypothesisContinue, left=of support] (continue) {继续或带条件继续\\冻结边界};
    \node[quantHypothesisRepair, right=of support] (stop) {停止\\如实保留不支持};
    \draw[quantArrow] (valid) -- node[above] {否} (repair);
    \draw[quantArrow] (valid) -- node[right] {是} (enough);
    \draw[quantArrow] (enough) -- node[above] {否} (insufficient);
    \draw[quantArrow] (enough) -- node[right] {是} (support);
    \draw[quantArrow] (support) -- node[above] {是} (continue);
    \draw[quantArrow] (support) -- node[above] {否} (stop);
  \end{tikzpicture}
  \caption{支持、返工、停止与证据不足的阶段决策。返工要求存在可定位的数据或实现问题。}
  \label{fig:ch01-decision-flow}
\end{figure}

\section{停止条件}

停止条件在实验前登记，但不填写未经确认的数值阈值。\cref{tab:ch01-stop-conditions} 给出本案例的原则级清单。

\begin{table}[htbp]
  \centering
  \footnotesize
  \caption{停止条件与阶段状态}
  \label{tab:ch01-stop-conditions}
  \begin{tabularx}{\textwidth}{p{5.5cm} p{3.0cm} Y}
    \toprule
    观察状态 & 阶段状态 & 解释 \\
    \midrule
    重大前视偏差；历史成分或关键数据版本无法恢复 & 返回修复；无法恢复时停止 & 证据链失去时间或版本基础 \\
    覆盖率不足以支持评价 & 证据不足 & 不得缩小分母制造“通过” \\
    关键参数只有孤立尖峰 & 停止或带条件继续 & 不把尖峰替换为主基线 \\
    方向在多数预注册分段不一致 & 停止或证据不足 & 区分稳定反证与样本不足 \\
    结果主要由无法解释的非目标暴露驱动 & 返回机制审查或停止 & 不能把未解释暴露当作原机制证据 \\
    换手或成本压力与研究目标明显冲突 & 停止或带条件继续 & 适用边界必须冻结并披露 \\
    测试集或最终留出集被反复用于调参 & 停止本轮确认性结论 & 新想法只能进入下一研究版本 \\
    因子缺乏增量信息 & 停止其进入正式组合 & 不因单因子表面相关而保留 \\
    数据和实现可信，但证据持续不支持假设 & 停止 & 不以返工名义继续搜索正结果 \\
    \bottomrule
  \end{tabularx}
\end{table}

具体项目可以把“继续、带条件继续、返回实现或数据修复、停止、证据不足”写成枚举字段。阈值需实际确认，但状态语义不能随结果改变。

\section{中期动量假设卡}

\begin{casebox}[title={中期动量假设卡}]
\textbf{因子名称与研究问题：}中期动量；历史中期相对强势是否预测下一周度持有期的相对收益。\textbf{候选机制：}信息扩散缓慢、反应不足或机构交易延续。\textbf{可计算代理与方向：}使用点时价格构造的中期累计收益，统一为数值越高、预期相对收益越高。

\textbf{数据与历史可得时间：}当时可得的交易日、价格、复权或公司行动信息、历史股票池；每个信号日只用截止时点以前记录。\textbf{主窗口：}约 12 个月历史并跳过最近约 1 个月，仅为教学基线，不是行业唯一标准或回测优选结果。\textbf{标签：}下一次调仓成交代理时点之间的未来相对收益。

\textbf{主要检验：}覆盖率、RankIC 方向、分组单调性、暴露、换手、成本与稳健性。\textbf{预注册敏感性：}少量邻近回看、跳过与最小历史观测设定。\textbf{可能暴露与交易风险：}高波动、流动性、规模、拥挤和较高换手。

\textbf{支持：}预注册证据方向一致、可解释且具备基本可交易性；\textbf{否定：}可信数据下方向持续不一致、主要由非目标暴露驱动或交易压力冲突；\textbf{证据不足：}历史或覆盖不足；\textbf{停止：}前视、复权错误无法修复、孤立窗口尖峰或反复窗口搜索。\textbf{主要失效机制：}市场快速反转、拥挤解除、流动性冲击，以及价格或复权处理错误。
\end{casebox}

\section{账面市值比价值假设卡}

\begin{casebox}[title={账面市值比价值假设卡}]
\textbf{因子名称与研究问题：}账面市值比价值；在点时财务信息下，相对较高的账面市值比是否预测下一周度持有期的相对收益。\textbf{候选机制：}市场对暂时不受欢迎或定价偏低公司的风险与现金流预期反应不充分。\textbf{可计算代理与方向：}可得权益除以同一口径市值，统一为数值越高、预期相对收益越高。

\textbf{数据与历史可得时间：}权益报告期、公告/可得时间、修订版本和信号日市值；财务数据只能在公告或保守可得时点后使用。\textbf{主口径：}权益定义、市值时点与单位需实际确认；负权益不得无说明纳入。\textbf{标签：}与统一研究任务相同。

\textbf{主要检验：}覆盖、方向、单调性、行业和市值暴露、换手、成本、稳健性与增量。\textbf{预注册敏感性：}少量保守财务滞后、市值口径或预先定义的变换。\textbf{可能暴露与交易风险：}行业资产结构、市值、流动性及困境公司暴露。

\textbf{支持：}点时口径稳定且多维证据一致；\textbf{否定：}可信数据下证据持续不支持或完全由非目标暴露解释；\textbf{证据不足：}公告历史、修订版本或覆盖不足；\textbf{停止：}历史可得时间无法恢复、负权益/单位错误或无界口径搜索。\textbf{主要失效机制：}价值陷阱、行业不可比、会计质量问题和估值持续重定价。
\end{casebox}

\section{ROE 类质量假设卡}

\begin{casebox}[title={ROE 类质量假设卡}]
\textbf{因子名称与研究问题：}ROE 类质量；点时可得的持续盈利能力是否预测下一周度持有期的相对收益。\textbf{候选机制：}盈利效率可能反映经营质量、资本配置或竞争优势，但市场未必即时完全定价。\textbf{可计算代理与方向：}所有组成项均点时可得的 TTM 利润与可比权益口径之比，数值越高、预期相对收益越高。

\textbf{数据与历史可得时间：}每个 TTM 组成项的报告期、公告/可得时间、修订版本和权益分母；不得用未来年度报告补齐历史指标。\textbf{主口径：}利润语义、平均或期末权益、小分母、负权益和一次性利润处理需实际确认。\textbf{标签：}与统一研究任务相同。

\textbf{主要检验：}覆盖、方向、单调性、行业、规模与杠杆暴露、换手、成本、稳健性和增量。\textbf{预注册敏感性：}少量 TTM 构造、可比权益与小分母规则。\textbf{可能暴露与交易风险：}行业不可比、杠杆抬高 ROE、低覆盖和高估值拥挤。

\textbf{支持：}点时组成完整且预注册证据一致；\textbf{否定：}证据持续不支持、质量已被估值充分反映或结果主要来自杠杆；\textbf{证据不足：}TTM 链或可比权益历史不足；\textbf{停止：}未来报告补齐、一次性利润污染无法识别或分母规则失控。\textbf{主要失效机制：}盈利均值回归、会计噪声、行业结构差异、杠杆和高质量已充分定价。
\end{casebox}

\section{统一研究任务说明}

三张因子卡必须挂在同一研究任务下，避免各自更换股票池、标签或判断规则。\cref{tab:ch01-research-task} 冻结贯穿案例的最小字段；“需实际确认”不是空白许可，而是待确认项的显式状态。

\begingroup
\scriptsize
\setlength{\tabcolsep}{3pt}
\begin{longtable}{@{}p{2.55cm}p{4.95cm}p{2.55cm}p{4.95cm}@{}}
\caption{统一研究任务冻结字段}\label{tab:ch01-research-task}\\
\toprule
字段 & 冻结教学设定 & 字段 & 冻结教学设定 \\
\midrule
\endfirsthead
\toprule
字段 & 冻结教学设定 & 字段 & 冻结教学设定 \\
\midrule
\endhead
\texttt{research\_id} & 唯一研究编号，不复用 & \texttt{research\_question} & 三个候选代理是否提供可解释、增量且基本可交易的横截面信息 \\
\texttt{market} & A 股普通股票日频横截面研究 & \texttt{universe\_definition} & 有历史成分记录的可替换宽基股票池；具体指数需实际确认 \\
\texttt{signal\_frequency} & 周频 & \texttt{rebalance\_frequency} & 周频 \\
\texttt{signal\_cutoff} & 每周最后一个交易日收盘后 & \texttt{execution\_timing} & 下一交易日；价格代理需实际确认 \\
\texttt{holding\_period} & 至下一调仓成交代理时点 & \texttt{label\_definition} & 两个成交代理时点之间未来收益；细节需实际确认 \\
\texttt{candidate\_factors} & 中期动量、账面市值比、ROE 类质量 & \texttt{main\_preprocessing} & 有效筛选、逐日 MAD、固定方向与逐日标准化；参数见第 \cref{ch:preprocessing} \\
\texttt{main\_evaluation\_metrics} & 覆盖、RankIC、单调性、暴露、换手、成本、稳健性与增量 & \texttt{robustness\_plan} & 时间顺序切分与少量预注册邻域；边界需实际确认 \\
\texttt{portfolio\_form} & 不加杠杆、不做空的约束型多头组合 & \texttt{cost\_scope} & 显性费用、滑点、冲击与不可成交；口径需实际确认 \\
\texttt{primary\_decision\_rule} & 多维预注册证据闸门，不以单一平均值决定 & \texttt{stop\_conditions} & \cref{tab:ch01-stop-conditions} 的原则清单与实际确认阈值 \\
\texttt{data\_requirements} & PIT 价格、财务公告与修订、历史成分、交易状态和公司行动 & \texttt{known\_limitations} & 数据版本、成交代理、费用、阈值和机构口径需实际确认 \\
\texttt{owner} & 研究与复核责任人需实际确认 & \texttt{version} & 冻结配置版本，任何实质变更递增 \\
\bottomrule
\end{longtable}
\endgroup

\section{实验登记和版本冻结}

\texttt{research\_id} 标识一项问题，\texttt{experiment\_id} 标识该问题下的一次具体运行。旧实验采用追加式记录；修复或改配置必须生成新编号，不得覆盖不利结果。

\begin{table}[htbp]
  \centering
  \footnotesize
  \caption{最小实验登记模板}
  \label{tab:ch01-experiment-registry}
  \begin{tabularx}{\textwidth}{p{3.0cm} Y p{3.0cm} Y}
    \toprule
    字段 & 记录内容 & 字段 & 记录内容 \\
    \midrule
    \texttt{experiment\_id} & 唯一运行编号 & \texttt{research\_id} & 所属研究任务 \\
    \texttt{hypothesis\_version} & 假设卡版本 & \texttt{config\_version} & 参数与规则版本 \\
    \texttt{data\_version} & 数据快照与可得规则 & \texttt{code\_commit} & 可复现代码提交 \\
    \texttt{sample\_role} & 训练、验证、测试或最终留出 & \texttt{reason\_before\_run} & 看到结果前的运行理由 \\
    \texttt{status} & 成功、失败或证据不足 & \texttt{result\_location} & 不覆盖的结构化结果位置 \\
    \texttt{decision} & 五类阶段状态之一 & \texttt{access\_count} & 测试与留出查看次数 \\
    \texttt{reviewer\_note} & 独立核验与异议 & \texttt{created\_at} & 登记时间，不替代数据时点 \\
    \bottomrule
  \end{tabularx}
\end{table}

\begin{practicebox}
\textbf{[正确实践] 自动核对清单：}研究与实验编号唯一；主配置在运行前冻结；因子、股票池、标签、数据和代码版本非空；信号截止早于标签起点；敏感性候选属于登记集合；探索结果不能写回主结论；测试与最终留出访问次数可追溯；失败实验不被删除；停止状态阻止进入下一阶段。
\end{practicebox}

\section{一个匿名研究决策案例}

\begin{casebox}
\textbf{仅用于解释研究设计和决策纪律，不代表真实证券、真实数据库或实证研究结果。}

某研究者提出“某个盈利能力指标可能预测未来相对收益”。错误流程是先尝试多个利润口径，再改变财务滞后和持有期，只保留最有利结果，最后补写经济逻辑。这样无法判断证据来自原想法，还是来自大量未披露选择。

正确流程是先登记盈利能力为何可能与未来相对收益相关，冻结点时可得的利润和权益代理、主标签与股票池，只列少量有理由的敏感性，并事前写清支持、不支持和证据不足。每个失败、报错或不利实验都保留编号；若看到结果后产生新利润口径，它只能成为下一版本的探索性假设，不能反向成为本轮确认性证据。
\end{casebox}

案例不包含证券、日期、阈值或指标数值，因为本节只演示决策纪律。即使最终结果不支持研究假设，只要数据与实现可信，正确动作也是记录并停止，而不是继续寻找正结果。

\section{YAML/Python 风格最小研究配置}

\cref{lst:ch01-minimal-config} 只展示配置结构。\texttt{ACTUAL\_CONFIRM} 表示在运行真实研究前必须由数据和研究负责人确认，不能原样当作参数。

\begin{lstlisting}[caption={最小研究配置（结构示例，非真实项目参数）},label={lst:ch01-minimal-config}]
research_config = {
    "research_id": "RESEARCH_ID_TO_ASSIGN",
    "universe": {
        "type": "historical_broad_index_membership",
        "index_id": "ACTUAL_CONFIRM",
    },
    "signal_schedule": {
        "frequency": "weekly",
        "cutoff": "last_trading_day_after_close",
    },
    "execution_proxy": "ACTUAL_CONFIRM_NEXT_TRADING_DAY_PROXY",
    "holding_period": "to_next_rebalance_execution_proxy",
    "factors": ["medium_term_momentum", "book_to_price", "roe_quality"],
    "primary_metrics": [
        "coverage", "rank_ic", "group_monotonicity", "exposures",
        "turnover", "cost_pressure", "robustness", "incremental_value",
    ],
    "sensitivity_plan": "PREREGISTERED_NEARBY_SETTINGS_ONLY",
    "stop_conditions": "VERSIONED_PRINCIPLE_LIST_AND_ACTUAL_GATES",
    "versions": {
        "hypothesis": "v1", "data": "ACTUAL_CONFIRM",
        "universe": "ACTUAL_CONFIRM", "label": "v1", "code": "COMMIT_ID",
    },
}
\end{lstlisting}

\section{常见错误}

\begin{itemize}
  \item \textbf{把机制当证据：}逻辑合理不能代替覆盖率、方向、稳健性和交易压力检查。
  \item \textbf{把公式当假设：}一个可运行表达式并未说明预测对象、否定条件或评价规则。
  \item \textbf{用未来信息修补代理：}财务报告、历史成分或成交状态越过信号截止，会使结论失去时间意义。
  \item \textbf{用单个平均指标决定去留：}它会掩盖方向翻转、非目标暴露、换手和成本。
  \item \textbf{把证据不足写成通过：}有效证券或有效期不足时，应返回状态而非放宽分母。
  \item \textbf{把不支持写成数据问题：}只有可定位错误才允许返工；可信的不利证据必须保留。
  \item \textbf{让探索结果改写主分析：}看到结果后产生的设定必须新建版本。
  \item \textbf{以正收益为结业要求：}这会诱导隐藏失败、调参和不诚实的范围扩张。
\end{itemize}

\section{本章产出}

\begin{outputbox}
\textbf{[本章产出]} 完成后应保存：一份带唯一 \texttt{research\_id} 和版本号的统一研究任务；三张因子假设卡；主分析与预注册敏感性清单；原则级停止条件；最小实验登记模板；待实际确认的股票池、数据、价格代理、成本、阈值与责任人清单。所有内容均在第一次查看结果前冻结。
\end{outputbox}

研究任务进入第 2 章的最低门槛不是“看起来会赚钱”，而是每个字段都有明确值或明确的“需实际确认”状态，且支持、不支持、返工和证据不足都能由事前规则产生。

\section{章节自测}

\begin{quizbox}
\textbf{[章节自测]}
\begin{enumerate}
  \item “低估值股票可能反弹”分别需要怎样的机制、代理、标签和否定证据？
  \item 为什么平均 RankIC 方向有利仍不足以支持 \(H_1\)？
  \item 主分析、预注册敏感性和探索性分析在何时定义，分别能否改变本轮主结论？
  \item 发现财务公告时间错误与发现可信结果不支持假设，为什么对应不同阶段状态？
  \item 下一交易日停牌为什么不能回头删除信号日股票？
  \item 中期动量、账面市值比和 ROE 各自最需要警惕哪些非目标暴露或数据错误？
  \item 什么情况下应写“证据不足”，而不是“支持”或“不支持”？
  \item 为什么测试集被反复用于调参后，不能通过改名恢复其样本外身份？
\end{enumerate}
\end{quizbox}

\section{与第 2 章《数据与时间边界》的衔接}

本章回答“要研究什么、什么证据会改变决定”；第 \cref{ch:data-timing} 接着回答“在每个历史时点究竟能知道什么”。之后，第 \cref{ch:universe-labels} 冻结历史股票池、可交易边界与收益标签，第 \cref{ch:factor-construction,ch:preprocessing} 实现代理并完成基础处理，第 \cref{ch:single-factor,ch:robustness} 形成单因子与稳健性证据。只有候选因子通过这些前向闸门，才进入第 \cref{ch:factor-combination,ch:portfolio,ch:backtest} 的合成、组合和含成本回测。

因此，下一章不会先读取结果，而会从字段级时间语义开始：观察日期、报告期、公告日期、数据可得日期、生效日期、修订版本与信号截止时间。若这些边界不能被恢复，本章写得再漂亮的假设也无法得到可信检验。

\chapter{数据与时间边界}

\section*{待编写}

% TODO: 按 P03 要求完整编写本章。

\chapter{股票池、可交易边界与收益标签}
\label{ch:universe-labels}

\section{本章导读与学习目标}

第 \cref{ch:data-timing} 已经建立点时信息集：站在任一历史时点，只允许使用当时已经可得的信息。本章在这一基础上继续回答两个更具体的问题：信号日的横截面究竟由哪些证券构成，以及每条信号记录应该对应哪一段未来收益。这里最重要的不是记住一份过滤清单，而是把\keyterm{样本资格、目标决策、成交结果和持仓状态}分开保存。

完成本章后，读者应能够：

\begin{enumerate}
  \item 用历史宽基成分和稳定证券标识重建信号日基础股票池；
  \item 判断一条规则应在信号截止前、目标组合生成时，还是下一交易日成交时生效；
  \item 区分研究股票池、信号有效股票池、目标组合候选、成交日可交易集合与实际持仓；
  \item 为每次纳入、剔除、未成交与持仓延续记录原因码；
  \item 生成周度信号日、成交代理时点和非重叠标签窗口；
  \item 区分标签缺失、订单无法成交、证券退市和实际持仓延续；
  \item 通过边界测试和人工穿行检查，证明没有用未来状态改写过去样本。
\end{enumerate}

\begin{principlebox}
\textbf{[通用原则]} 证券是否属于信号日股票池，只能由信号截止及以前的信息决定。下一交易日才出现的停牌、价格限制、成交价格缺失等状态，只能在成交模拟阶段处理；它们不能反向删除已经生成的信号日记录。
\end{principlebox}

\section{横截面样本为什么不仅是一张股票列表}

一张只含\texttt{signal\_date}和\texttt{security\_id}的列表，无法说明证券为何出现、何时被排除，也无法区分“研究时存在但后来无法买入”和“研究时本就不属于样本”。如果直接把最终成交证券当作因子研究样本，下一交易日的状态就会倒流到信号日，样本分布也随之被未来信息改变。

可审计的横截面底表应至少同时保留三类信息。第一类是\keyterm{成员证据}，说明证券在信号日是否属于历史基础池；第二类是\keyterm{信号日资格证据}，说明只利用信号截止前信息应用过滤后的结果；第三类是\keyterm{后续事件证据}，记录目标权重、订单、成交、未成交原因和实际持仓。它们可以通过稳定主键关联，却不应压缩成一个会被后续状态覆盖的布尔列。

\begin{teachingbox}
\textbf{[教学简化]} 贯穿案例继续采用“宽基股票池三因子周频约束多头策略”：每周最后一个交易日收盘后生成信号，下一交易日使用可配置的日频价格代理模拟成交。本章只讨论集合、时点和标签，不预设真实指数、价格字段或交易制度细节。
\end{teachingbox}

\section{研究股票池、信号有效股票池、目标组合候选、成交日可交易集合和实际持仓的区别}

令 \(t\) 表示某次周度信号日，令 \(\tau(t)\) 表示与该信号对应的下一交易日成交代理时点。研究流程中至少存在以下五个集合：

\begin{itemize}
  \item \(\mathcal{U}^{\mathrm{research}}_t\)：信号日的历史研究股票池，以当时有效的宽基成分为起点；
  \item \(\mathcal{U}^{\mathrm{signal}}_t\)：仅用信号截止及以前信息过滤后，允许计算和比较信号的证券；
  \item \(\mathcal{C}_t\)：经因子得分、排序和组合约束后进入目标组合候选的证券；
  \item \(\mathcal{X}_{\tau(t)}\)：到成交阶段依据当时交易状态与成交代理可用性判断的可执行证券；
  \item \(\mathcal{H}_{\tau(t)}\)：订单执行后真正持有的证券，包括本次成交形成的持仓以及因无法卖出而延续的旧持仓。
\end{itemize}

\cref{fig:ch03-universe-hierarchy} 展示了五个集合的职责边界。前两个集合决定研究横截面，目标候选表达“想持有什么”，成交日集合表达“此刻能做什么”，实际持仓则表达“最终发生了什么”。它们通常相互关联，但不能假定逐层都是简单子集：例如原持仓无法卖出时，它可能不再属于本期目标候选，却仍然留在实际持仓中。

\begin{figure}[htbp]
  \centering
  \begin{tikzpicture}[
    node distance=6mm,
    every node/.style={font=\small},
    quantUniverseFlow/.style={quantBlueCard, minimum width=7.2cm, align=center},
    quantUniverseNote/.style={quantCard, text width=5.5cm, align=left}
  ]
    \node[quantUniverseFlow] (research) {历史研究股票池\\\(\mathcal{U}^{\mathrm{research}}_t\)};
    \node[quantUniverseFlow, below=of research] (signal) {信号有效股票池\\\(\mathcal{U}^{\mathrm{signal}}_t\)};
    \node[quantUniverseFlow, below=of signal] (candidate) {目标组合候选\\\(\mathcal{C}_t\)};
    \node[quantOrangeCard, minimum width=7.2cm, align=center, below=of candidate]
      (execution) {成交日可执行集合\\\(\mathcal{X}_{\tau(t)}\)};
    \node[quantGreenCard, minimum width=7.2cm, align=center, below=of execution]
      (holding) {订单处理后的实际持仓\\\(\mathcal{H}_{\tau(t)}\)};

    \node[quantUniverseNote, right=7mm of research] {历史成分与信号日前已生效主数据};
    \node[quantUniverseNote, right=7mm of signal] {只用 \(\mathcal{I}_t\) 过滤，保留原因码};
    \node[quantUniverseNote, right=7mm of candidate] {因子排序与组合约束给出目标};
    \node[quantUniverseNote, right=7mm of execution] {此处才读取下一交易日交易状态};
    \node[quantUniverseNote, right=7mm of holding] {成交、未成交与旧持仓延续共同决定};

    \draw[quantArrow] (research) -- (signal);
    \draw[quantArrow] (signal) -- (candidate);
    \draw[quantArrow] (candidate) -- (execution);
    \draw[quantArrow] (execution) -- (holding);
  \end{tikzpicture}
  \caption{研究池、信号池、目标组合、成交与实际持仓的层级关系}
  \label{fig:ch03-universe-hierarchy}
\end{figure}

\section{使用历史宽基成分构建信号日基础股票池}

贯穿案例的基础池来自可配置宽基指数的\keyterm{历史成分记录}。对每个信号日 \(t\)，应选择在该日已经生效、且有效区间覆盖 \(t\) 的成员记录，形成 \(\mathcal{U}^{\mathrm{research}}_t\)。成员表需要保留来源版本、生效区间和匹配依据，使任一历史横截面能够重放。

这里不能用今天的指数成分覆盖过去，也不能先取今天仍存续的证券再向历史回填。前者会遗漏历史上被调出的证券，后者会遗漏后来退市或代码变化的证券，两者都会造成幸存者偏差。第 \cref{ch:data-timing} 已解释点时连接与生效区间，本节只强调其集合含义：\keyterm{基础池是按信号日逐期生成的历史状态，而不是一张静态名单。}

\begin{warningbox}
用当前成分或当前存续股票列表重建历史，会把“后来留下来”误写成“当时就应被研究”。即使因子值和收益标签都没有显式引用未来日期，这种样本选择本身也已经使用了未来信息。
\end{warningbox}

\section{稳定证券标识、上市、退市和代码变更}

证券代码是展示标识，不一定适合作为跨期唯一主键。研究底表应使用能够跨代码变更保持一致、同时区分不同证券实体的稳定标识，再通过带生效区间的映射表获取某日代码。这样，代码变化不会把一条证券历史拆成两只“新证券”，代码复用也不会错误拼接不同实体。

上市和退市同样应作为有时间边界的主数据事件。信号日之前尚未上市的证券不属于当期样本；信号日仍有效、后来退市的证券不能因为最终结局而从历史中消失。退市后的估值、最后可交易时点、现金或其他公司行动如何进入标签与持仓账本，需要另行确定，但原始成员记录和信号记录必须保留。

\begin{confirmbox}
\textbf{[需实际确认]} 应核验数据源的稳定证券主键、代码变更映射、生效时间精度、上市与退市字段语义，以及公司行动后的价格和持仓调整方式。[机构实现] 这些字段在生产表中的名称、版本和关联键也必须由团队数据字典给出，本章不作推断。
\end{confirmbox}

\section{信号日前可得的过滤规则}

过滤规则的判断标准不是“它看起来像交易限制”，而是\keyterm{该条件在什么时点可知，以及它改变哪个对象}。若某项状态在信号截止前已经生效，并且研究方案预先约定把它排除在因子比较范围之外，它可以改变 \(\mathcal{U}^{\mathrm{signal}}_t\)。若状态直到 \(\tau(t)\) 才能观察到，它只能改变订单是否成交，不能改写 \(\mathcal{U}^{\mathrm{signal}}_t\)。

\cref{tab:ch03-filter-timing} 给出规则归属。表中“可使用”是时间约束，并不代表一定应该过滤；是否启用、阈值多少以及异常如何处理，仍由研究协议决定。

{
\setlength{\tabcolsep}{3.2pt}
\small
\begin{longtable}{p{2.35cm} p{2.75cm} p{2.65cm} p{2.8cm} p{3.55cm}}
\caption{过滤规则及其生效时点}\label{tab:ch03-filter-timing}\\
\toprule
条件类别 & 信号日可得证据 & 应影响的阶段 & 不应执行的操作 & 口径状态 \\
\midrule
\endfirsthead
\toprule
条件类别 & 信号日可得证据 & 应影响的阶段 & 不应执行的操作 & 口径状态 \\
\midrule
\endhead
历史宽基成员 & 覆盖 \(t\) 的已生效历史成分 & 研究股票池 & 用当前成分回填历史 & [通用原则] \\
上市状态 & 截止 \(t\) 已生效的证券主数据 & 研究池或信号池 & 用未来退市结果删除过去 & 字段语义需确认 \\
上市时间长度 & 截止 \(t\) 可计算的历史长度 & 信号池 & 把未确认阈值写成统一规则 & 阈值需确认 \\
特殊处理状态 & 截止 \(t\) 已生效且当时可知的状态 & 信号池（若协议启用） & 使用后来补记状态提前过滤 & 状态与生效时间需确认 \\
历史停牌状态 & 截止 \(t\) 已知的历史记录 & 信号池（若协议启用） & 把 \(\tau(t)\) 状态提前到 \(t\) & 过滤目的需确认 \\
历史流动性 & 以 \(t\) 或更早结束的历史窗口 & 信号池或目标约束 & 使用未来成交量计算窗口 & 窗口与阈值需确认 \\
成交日停牌 & 到 \(\tau(t)\) 才确定的状态 & 成交模拟 & 回删信号日样本或因子值 & [通用原则] \\
成交日价格限制 & 到 \(\tau(t)\) 才确定的状态 & 成交模拟 & 回删目标候选与信号记录 & 规则与字段需确认 \\
成交价格代理缺失 & 成交阶段观察到的价格可用性 & 成交与标签状态 & 把缺失解释为历史不存在 & 代理与缺失语义需确认 \\
\bottomrule
\end{longtable}
}

\section{上市时间、特殊处理状态、停牌、价格限制和流动性条件的时点问题}

上市时间长度可以由信号日前已知的上市日期计算，因此从时间上可以作为信号日过滤；但最低历史长度并非天然常数。特殊处理状态也只有在其状态及生效时间于信号截止前可知时，才可能进入信号池过滤。是否过滤以及如何处理状态变更，是研究口径，而不是本章替团队做出的制度判断。

停牌需要区分两个问题：信号日以前的历史停牌记录可能影响因子覆盖或预设资格；下一交易日是否停牌则只决定订单能否执行。价格限制同理，成交日的状态不能反向参与信号样本选择。流动性条件还要区分以历史窗口度量的候选资格与成交日实际容量：前者的窗口必须在 \(t\) 截止，后者属于执行模拟。

\begin{confirmbox}
\textbf{[需实际确认]} 上市时间阈值、特殊处理状态的定义与生效时点、停复牌记录语义、价格限制判断、历史流动性窗口、成交量约束及具体成交价格代理，都依赖研究时期、交易所规则、数据供应商和团队协议。实施前应逐项核验原始字段、时间戳、缺失含义与历史制度版本。
\end{confirmbox}

\section{为什么不能用下一交易日状态反向删除信号日证券}

假设匿名证券在信号日 \(t\) 属于历史成分，且使用 \(\mathcal{I}_t\) 计算出的因子数据完整。它因此合法地属于 \(\mathcal{U}^{\mathrm{signal}}_t\)。若到 \(\tau(t)\) 才发现无法成交，正确结果是目标订单具有“未成交”状态，而不是让这只证券从信号日横截面消失。

回删会产生两类问题。其一，研究样本隐含知道了下一交易日状态，形成前视选择；其二，因子分位数、排序阈值和目标权重可能被重新计算，使其他证券的信号也受到未来信息影响。即使被删证券本身没有收益标签，横截面统计仍然可能被系统性美化。

\begin{principlebox}
信号快照应当不可变。成交模拟可以新增 \texttt{execution\_status}、\texttt{unfilled\_reason} 和实际持仓记录，却不得覆盖信号日的 \texttt{in\_research\_universe}、\texttt{signal\_eligible} 或原始因子值。
\end{principlebox}

\section{纳入与剔除原因码}

单一的 \texttt{eligible=True/False} 无法支持复核。每个阶段应同时记录结果、原因码、判断时点和规则版本。原因码最好采用稳定机器码，文字说明可以随文档更新；多条规则同时命中时，应保留完整原因集合或明确优先级，而不是只留下最后一次覆盖的结果。

\cref{tab:ch03-reason-codes} 是教学用最小字典。它展示证据链的结构，不代表任何机构现有字段或最终枚举。

\begin{table}[htbp]
  \centering
  \small
  \caption{纳入、剔除、成交与持仓原因码示例}
  \label{tab:ch03-reason-codes}
  \begin{tabularx}{\textwidth}{p{3.35cm} p{2.6cm} Y Y}
    \toprule
    教学原因码 & 判断阶段 & 含义 & 必须保留的证据 \\
    \midrule
    \texttt{BASE\_MEMBER} & 研究池 & 信号日属于历史基础池 & 成分生效区间与版本 \\
    \texttt{EXCL\_NOT\_MEMBER} & 研究池 & 信号日不属于历史基础池 & 未匹配或区间不覆盖 \\
    \texttt{SIGNAL\_ELIGIBLE} & 信号池 & 信号日前规则全部通过 & 逐项规则结果与版本 \\
    \texttt{EXCL\_SIGNAL\_STATUS} & 信号池 & 信号日前某项资格未通过 & 状态值、可得日与规则 \\
    \texttt{EXCL\_INPUT\_MISSING} & 信号池 & 计算信号所需输入缺失 & 缺失字段类别与检查时间 \\
    \texttt{TARGET\_SELECTED} & 目标组合 & 经排序和约束进入目标 & 分数、目标权重与配置 \\
    \texttt{ORDER\_UNFILLED} & 成交阶段 & 订单在成交代理时点未执行 & 订单、当时状态与失败原因 \\
    \texttt{HOLD\_CARRY} & 持仓阶段 & 旧持仓因未完成卖出而延续 & 前期持仓、退出订单与状态 \\
    \texttt{LABEL\_MISSING} & 标签阶段 & 约定窗口的价格代理不完整 & 缺失端点及原始可用性 \\
    \bottomrule
  \end{tabularx}
\end{table}

\section{周度信号日和下一交易日的生成}

周度日程必须从历史交易日历生成，而不是用自然日加减。对每个研究周，先在日历中选择最后一个交易日作为 \(t\)，再从同一日历中查找严格晚于 \(t\) 的首个交易日，作为与本期信号对应的成交日。节假日会自然改变二者间的自然日距离，但不会改变“下一交易日”的定义。

生成日程后，应把 \texttt{signal\_date}、\texttt{execution\_date}、日历版本和规则版本固化为调仓日程表。若研究期末没有下一交易日，最后一条信号可以保留，但其成交与标签状态必须明确标记为不可观察，不能凭自然日补齐。

\begin{practicebox}
正确顺序是：先从历史交易日历生成全部周度信号日，再映射各自下一交易日，最后把日程表连接到成员、信号、订单和标签。人工抽查至少覆盖跨周节假日、研究期起止和日历异常日期。
\end{practicebox}

\section{成交代理、持有期和未来收益标签}

因子预测目标必须与可执行时序一致。信号使用 \(t\) 日收盘后可得信息时，不能同时假设无条件以 \(t\) 日收盘价成交。本章用 \(\tau(t)\) 表示下一交易日的可配置成交代理时点；下一次调仓信号记作 \(t+1\)，其成交代理时点记作 \(\tau(t+1)\)。这里的 \(t+1\) 表示\keyterm{下一次预定周度信号}，不是下一个自然日。

令 \(P_i^{\mathrm{proxy}}(s)\) 为证券 \(i\) 在时点 \(s\) 按统一口径取得的成交价格代理，则本章的非重叠周度未来收益标签定义为

\begin{equation}
  y_{i,t}
  = \frac{P_i^{\mathrm{proxy}}\!\left(\tau(t+1)\right)}
  {P_i^{\mathrm{proxy}}\!\left(\tau(t)\right)} - 1 .
  \label{eq:ch03-weekly-label}
\end{equation}

\cref{eq:ch03-weekly-label} 的两个端点必须采用一致且可解释的代理口径；任何复权、公司行动与费用处理都要另外写入配置和版本记录。标签描述的是模型要预测的未来区间，不自动等同于实际组合收益。实际组合收益还取决于订单是否成交、成交数量、旧持仓延续、现金和费用等状态。

\begin{confirmbox}
\textbf{[需实际确认]} \(P^{\mathrm{proxy}}\) 最终采用何种价格、端点如何处理复权与公司行动、订单可成交比例、费用及研究标签是否要求真实可成交，都必须按数据源和团队研究协议确认。本章不把任何具体价格字段或费用口径写成行业统一事实。
\end{confirmbox}

\section{非重叠周度标签的定义}

按连续两次调仓的成交代理时点切分持有期，可将第 \(t\) 期标签窗口写成 \([\tau(t),\tau(t+1))\)。下一期从 \(\tau(t+1)\) 开始，因此相邻窗口共享边界但不重复计算同一持有区间。这种定义使周度信号、目标组合和未来标签在时间上对齐，也便于把每条标签追溯到两条日程记录。

\cref{fig:ch03-label-timeline} 中，成员资格与信号值冻结在 \(t\) 收盘截止；成交状态到 \(\tau(t)\) 才读取；标签覆盖两次成交代理时点之间的区间。若证券在任一端点缺少约定代理，标签应进入显式缺失状态，而不是擅自延长或缩短持有期。

\begin{figure}[htbp]
  \centering
  \begin{tikzpicture}[
    x=1cm,
    y=1cm,
    every node/.style={font=\small},
    quantLabelEvent/.style={quantBlueCard, text width=2.7cm, align=center},
    quantLabelExecution/.style={quantOrangeCard, text width=2.8cm, align=center}
  ]
    \draw[line width=1.2pt, draw=SoftBlue] (0,0) -- (14.8,0);
    \node[quantLabelEvent] (s0) at (1.1,1.25) {信号日 \(t\) 收盘截止\\冻结成员与信号};
    \node[quantLabelExecution] (e0) at (5.0,-1.25) {成交代理 \(\tau(t)\)\\读取当时交易状态};
    \node[quantLabelEvent] (s1) at (10.0,1.25) {下一信号日 \(t+1\)\\生成新目标};
    \node[quantLabelExecution] (e1) at (13.7,-1.25) {成交代理 \(\tau(t+1)\)\\标签窗口结束};
    \foreach \x in {1.1,5.0,10.0,13.7}
      \draw[draw=SoftBlue, line width=1pt] (\x,-0.16) -- (\x,0.16);
    \draw[quantArrow, draw=SoftGreen] (5.0,0.35) -- node[above, align=center]
      {持有期与标签窗口\\\([\tau(t),\tau(t+1))\)} (13.7,0.35);
    \draw[quantSoftArrow] (s0.south) -- (1.1,0.16);
    \draw[quantSoftArrow] (e0.north) -- (5.0,-0.16);
    \draw[quantSoftArrow] (s1.south) -- (10.0,0.16);
    \draw[quantSoftArrow] (e1.north) -- (13.7,-0.16);
  \end{tikzpicture}
  \caption{周度信号、成交代理与非重叠标签窗口}
  \label{fig:ch03-label-timeline}
\end{figure}

\section{缺失标签、无法成交和退市样本的区别}

\keyterm{标签缺失}表示无法按约定端点计算 \cref{eq:ch03-weekly-label}；\keyterm{无法成交}表示订单在成交阶段没有按规则执行；\keyterm{退市样本}表示证券实体经历了有时间边界的主数据事件。这三种状态可能同时出现，但语义不同，不能用同一个空值替代。

证券有合法信号记录、却在 \(\tau(t)\) 无法买入时，信号横截面记录仍存在，订单日志记录未成交，实际持仓可能保持为零。若原本持有的证券无法卖出，则目标权重可以为零，但实际持仓必须延续并留痕。若标签端点价格缺失，应保留信号与缺失原因；不能把“没有标签”解释成“这只证券在历史上不存在”。

退市、长期停牌和公司行动尤其需要保留状态链。它们可能影响代理价格、持仓估值、现金流和退出方式，但具体处理依赖数据源和研究口径。若无法可靠恢复，应把样本状态标为待确认或不可评估，并披露影响，而不是静默删除。

\begin{confirmbox}
\textbf{[需实际确认]} 退市证券的最后价格、现金或其他公司行动，长期停牌期间的估值，标签端点缺失后的处理，以及未成交订单的有效期，都需要结合历史制度、数据覆盖和团队状态机核验。任何默认填充值、延长窗口或强制退出规则都必须显式配置。
\end{confirmbox}

\section{一个极小教学案例}

以下匿名案例只用于解释集合与留痕逻辑，不代表真实证券、真实数据库或实证研究结果。证券“样本 A”在信号日属于历史宽基成分，所有信号日前过滤均通过，并被目标组合选中；到下一交易日成交阶段，系统才观察到该证券无法按约定代理成交。

\begin{casebox}
正确记录方式是保留“样本 A”在 \(t\) 日的成员、因子值、排名和目标订单；在 \(\tau(t)\) 的订单日志新增 \texttt{ORDER\_UNFILLED} 及待核验原因；若此前未持有，则实际持仓仍为零，若此前已有持仓且退出订单也无法执行，则以 \texttt{HOLD\_CARRY} 记录延续。不得重新计算 \(t\) 日横截面，也不得删除该证券的信号行。
\end{casebox}

\begin{table}[htbp]
  \centering
  \small
  \caption{匿名证券“样本 A”的事件留痕}
  \label{tab:ch03-anonymous-case}
  \begin{tabularx}{\textwidth}{p{2.6cm} p{2.9cm} Y Y}
    \toprule
    阶段 & 记录状态 & 正确动作 & 禁止动作 \\
    \midrule
    信号日基础池 & \texttt{BASE\_MEMBER} & 保留历史成分证据 & 因后续未成交删除成员行 \\
    信号日过滤 & \texttt{SIGNAL\_ELIGIBLE} & 冻结资格、因子值与排序 & 读取下一交易日状态重排 \\
    目标组合 & \texttt{TARGET\_SELECTED} & 生成带版本的目标订单 & 把目标与实际持仓混为一列 \\
    下一交易日 & \texttt{ORDER\_UNFILLED} & 记录状态、原因和未成交数量 & 回写 \texttt{signal\_eligible=False} \\
    实际持仓 & 零持仓或 \texttt{HOLD\_CARRY} & 延续真实账本并记录差异 & 假定目标权重已经实现 \\
    标签 & 有值或 \texttt{LABEL\_MISSING} & 按端点可用性独立判断 & 用零收益替代未知标签 \\
    \bottomrule
  \end{tabularx}
\end{table}

\section{pandas 风格最小代码或伪代码}

下面的代码强调数据流，而不是给出可直接用于生产的字段规范。历史成员、信号过滤、目标订单、成交状态和标签被保存在不同表中；成交结果通过稳定主键和日程键回连，但不会覆盖信号快照。

\begin{lstlisting}[
  caption={股票池、成交状态和周度标签的最小数据流},
  label={lst:ch03-universe-labels}
]
# schedule comes from a versioned historical trading calendar
schedule = weekly_signals.merge(next_trading_dates, on="signal_date")

# Membership and filters use only information available by signal_date.
signal = historical_membership.merge(schedule, on="signal_date")
signal["signal_eligible"] = (
    signal["is_member_asof"]
    & signal["listed_asof"]
    & signal["signal_inputs_ready"]
    & signal["signal_status_ok"]
)
signal_snapshot = signal.copy()  # immutable after signal cutoff

# Execution state is joined only after targets have been generated.
orders = targets.merge(
    execution_state,
    on=["security_id", "execution_date"],
    how="left",
)
orders["filled"] = (
    orders["execution_allowed"]
    & orders["entry_proxy"].notna()
)

# Label endpoints use one consistent, configurable price proxy.
labels = entry_proxy.merge(
    next_entry_proxy,
    on=["security_id", "signal_date"],
    how="left",
)
labels["future_return"] = (
    labels["next_proxy"] / labels["entry_proxy"] - 1
)
labels["label_missing"] = labels[
    ["entry_proxy", "next_proxy"]
].isna().any(axis=1)

assert signal_snapshot.equals(signal)
assert orders["execution_date"].gt(orders["signal_date"]).all()
\end{lstlisting}

\cref{lst:ch03-universe-labels} 中的 \texttt{execution\_allowed} 不是预设行业字段，而是成交状态机在 \(\tau(t)\) 根据经确认规则产生的结果。实际实现还应输出逐项原因码、配置版本和持仓账本，不能仅保留最终布尔值。

\section{边界测试和人工核对}

自动化测试应围绕集合是否在正确时点发生变化，而不仅是检查行数。建议至少覆盖以下边界：

\begin{enumerate}
  \item \textbf{节假日：}周末或节假日后的成交日必须来自下一历史交易日，不得用 \(t+1\) 个自然日代替；
  \item \textbf{成分进出：}生效日前不纳入，生效日及以后按历史记录判断；后来调出不得改写此前成员资格；
  \item \textbf{上市与代码变更：}上市前不得进入，代码变化前后通过稳定标识连续，且不同实体不得因代码复用误拼；
  \item \textbf{退市：}退市前的信号行保持存在，退市后的标签与持仓状态按原因码进入待确认流程；
  \item \textbf{停牌：}信号日通过、下一交易日才无法成交的证券仍留在信号快照，订单标记未成交；
  \item \textbf{价格限制或代理缺失：}只改变成交或标签状态，不改变历史成员、信号资格和原始排名；
  \item \textbf{标签缺失：}任一端点缺失时不默认填零，并能指出缺失端点、证券和对应窗口；
  \item \textbf{持仓延续：}目标退出但卖单未成交时，目标组合与实际持仓的差异必须可追溯。
\end{enumerate}

\begin{practicebox}
人工核对应抽取若干信号日和匿名证券，沿“历史成分记录 \(\rightarrow\) 信号过滤证据 \(\rightarrow\) 目标订单 \(\rightarrow\) 成交状态 \(\rightarrow\) 实际持仓 \(\rightarrow\) 标签端点”逐步穿行。至少选取一次跨节假日、一次成分变更、一次上市或退市、一次成交受阻和一次标签缺失，并由核验人记录预期状态、实际状态、差异原因与证据位置。
\end{practicebox}

\section{常见错误}

\begin{warningbox}
常见错误包括：用当前指数成分或当前存续证券回填历史；把研究池、目标组合和实际持仓压成同一个集合；用 \(\tau(t)\) 的停牌或价格限制回删 \(t\) 日样本；以自然日生成下一交易日；把目标权重当作已成交持仓；把缺失标签填成零收益；在标签两端混用不同价格口径；静默删除退市或长期停牌证券；以及只保存最终布尔值而不保存原因码和规则版本。
\end{warningbox}

发现这些问题时，不应只修正最终收益序列。应从最早受污染的集合快照开始重建，并比较修复前后的成员数、原因码分布、目标订单、成交状态和标签覆盖，确认未来状态不再影响过去横截面。

\section{本章产出}

\begin{outputbox}
本章结束后应提交：版本化的周度信号与成交日程表；逐信号日历史研究池和信号有效池；纳入与剔除原因码字典；目标订单、成交结果与实际持仓的关联键设计；\cref{eq:ch03-weekly-label} 对应的标签定义和端点状态；自动化边界测试结果；以及包含节假日、成分变化、上市或退市、停牌和标签缺失的人工穿行记录。
\end{outputbox}

这些产出描述的是可复核的数据接口，不包含任何策略收益、因子有效性或实证结论。在没有真实金融数据和经确认口径之前，不应生成结果数字。

\section{章节自测}

\begin{quizbox}
\begin{enumerate}
  \item 为什么下一交易日无法买入，不能成为删除信号日证券的理由？
  \item \(\mathcal{U}^{\mathrm{signal}}_t\)、\(\mathcal{C}_t\) 与 \(\mathcal{H}_{\tau(t)}\) 分别回答什么问题？
  \item 某证券目标权重为零但仍在实际持仓中，可能发生了什么？应保留哪些记录？
  \item 为什么“标签缺失”和“证券在历史中不存在”不能使用同一个状态？
  \item 如何用历史交易日历构造 \(t\)、\(\tau(t)\) 与 \(\tau(t+1)\)？
  \item \cref{eq:ch03-weekly-label} 的两个价格端点需要满足哪些一致性要求？
  \item 若只能获得当前指数成分，是否可以重建历史基础池？应如何披露限制？
  \item 请为“信号日有效、成交日未成交、旧持仓延续”设计一条完整证据链。
\end{enumerate}
\end{quizbox}

\section{与第 4 章《因子构造与实现核对》的衔接}

本章把每个信号日可参与横截面比较的证券集合冻结下来，并为后续结果定义了统一预测窗口。第 4 章将在 \(\mathcal{U}^{\mathrm{signal}}_t\) 上实现中期动量、账面市值比价值和 ROE 类质量三个教学因子，重点核对输入窗口、方向、缺失值与单证券计算路径。

衔接时必须继续保持两条边界：因子只能读取信号截止前可得数据，因子构造不能读取 \cref{eq:ch03-weekly-label} 的未来端点；成交日状态也不得反向参与因子缺失处理或横截面排序。这样，第 4 章得到的因子表才能与本章的成员快照和标签表通过稳定主键安全连接。

\chapter{因子构造与实现核对}
\label{ch:factor-construction}

\section{本章导读与学习目标}

前两章已经冻结了信号日点时信息集与可参与比较的证券集合：时间可得性见 \cref{ch:data-timing}，股票池和未来收益标签边界见 \cref{ch:universe-labels}。本章只在这些已经冻结的输入上完成三个教学因子的\keyterm{原始构造与实现核对}：中期动量、账面市值比价值和 ROE 类质量。这里的目标不是证明因子有效，而是证明“公式写的”和“程序算的”是同一件事。

完成本章后，读者应能够：

\begin{enumerate}
  \item 把经济假设写成公式、数据依赖、时间窗口、方向和有效条件；
  \item 区分沿单只证券历史展开的时间序列运算与同一信号日的横截面运算；
  \item 保证每个输入在信号截止前已经可得，且不访问未来标签或成交状态；
  \item 同时输出原始因子值、有效性标记、无效原因和输入最新可得日；
  \item 用单元测试与匿名证券人工手算核对实现；
  \item 用统一 schema、定义卡和版本号建立可复用、可审计的因子接口。
\end{enumerate}

\begin{principlebox}
[通用原则] 因子实现正确性的最低标准不是“结果看起来合理”，而是任一 \texttt{security\_id} 与 \texttt{signal\_date} 的输出都能追溯到已冻结输入、参数、公式、有效性判断和代码版本。未来收益不得用于选窗口、改方向或修补缺失值。
\end{principlebox}

本章不进行去极值、缩尾、排名、标准化、行业或市值中性化，也不计算 IC、分组收益或组合回测。这些步骤分别属于第 5 章之后的预处理、评价和组合流程。

\section{从经济假设到可计算因子}

经济假设通常以自然语言开始，例如“过去一段较长时期表现相对强的证券可能延续这一状态”“账面权益相对市值较高的证券可能具有更高价值暴露”“持续盈利且权益基础稳健的公司可能具有更高质量”。这些都是待验证假设，不是已经成立的事实。要把它们变成可计算因子，至少要回答五个问题：测量对象是什么，使用哪些输入，输入何时可得，参数如何预先固定，以及什么情况下拒绝给出数值。

一个完整转换过程应先登记研究假设和预期方向，再锁定点时数据、参数与有效条件，然后计算原始值，最后输出质量状态。若先看未来表现再决定窗口或翻转方向，程序即使没有显式连接标签表，也已经把未来结果带入了因子定义。

\begin{figure}[htbp]
  \centering
  \begin{tikzpicture}[
    node distance=3mm,
    every node/.style={font=\small},
    step/.style={bluecard, text width=2.25cm, minimum height=1.35cm, align=center},
    check/.style={greencard, text width=2.35cm, minimum height=1.35cm, align=center}
  ]
    \node[step] (hypothesis) {经济假设\\与预期方向};
    \node[step, right=of hypothesis] (data) {点时数据\\参数与窗口};
    \node[step, right=of data] (raw) {原始因子值\\\texttt{raw\_value}};
    \node[check, right=of raw] (valid) {有效性检查\\原因与版本};
    \node[step, right=of valid] (next) {第 5 章\\后续预处理};

    \draw[myarrow] (hypothesis) -- (data);
    \draw[myarrow] (data) -- (raw);
    \draw[myarrow] (raw) -- (valid);
    \draw[myarrow] (valid) -- (next);
  \end{tikzpicture}
  \caption{从研究假设到原始因子输出的流程（结构示意，非实证结果）}
  \label{fig:ch04-factor-pipeline}
\end{figure}

\cref{fig:ch04-factor-pipeline} 中“有效性检查”不是事后清洗。它与公式一起定义原始因子：历史不足、点时记录不可得或分母不合法时，程序应拒绝生成看似精确的数值，而不是等到下一章再处理。

\section{原始因子、处理后因子和交易信号的区别}

\keyterm{原始因子}是按定义卡公式直接计算、尚未经过第 5 章截面处理的数值，例如一个价格比值减一或两个点时会计量的比值。\keyterm{处理后因子}是在原始值及其有效状态基础上，经预先登记的缩尾、变换、标准化或暴露控制得到的可比较数值。\keyterm{交易信号}还要结合因子方向、多因子合成、股票池、组合约束与持仓状态，才能转化为目标权重或订单。

三者不能使用同一列反复覆盖。否则研究者无法判断一个极端值来自原公式、预处理还是组合约束，也无法重现某次版本变化。本章的 \texttt{raw\_value} 只保存原始值；其数值大小不代表已可直接比较、交易或用于组合优化。

\begin{warningbox}
把截面排名写回原始因子列，会永久丢失公式输出；把缺失值填零后再排名，会把“未知”伪装成“中性”；把目标持仓称为因子，则会混淆预测、组合和执行三个层次。
\end{warningbox}

\section{时间序列运算与横截面运算}

时间序列运算沿同一证券的日期轴移动，例如计算历史价格端点、滞后值和滚动观测数；横截面运算则在同一信号日比较不同证券，例如排名、标准化或行业内变换。动量的端点选择属于时间序列运算，账面市值比和 ROE 的点时比值也是逐证券计算。本章尚不进行横截面排名或标准化。

\cref{tab:ch04-ts-cs-comparison} 给出两类运算的防错边界。最危险的实现错误是排序或分组键不完整：未按证券分组的 \texttt{shift} 会把上一只证券的历史接到下一只证券，未按信号日分组的排名则会把不同日期混在一起。

\begin{table}[htbp]
  \centering
  \small
  \caption{时间序列运算与横截面运算对比}
  \label{tab:ch04-ts-cs-comparison}
  \begin{tabularx}{\textwidth}{p{2.6cm} Y Y Y}
    \toprule
    比较项 & 时间序列运算 & 横截面运算 & 本章边界 \\
    \midrule
    核心问题 & 同一证券过去发生了什么 & 同一时点不同证券如何比较 & 只输出逐证券原始值 \\
    分组键 & \texttt{security\_id} & \texttt{signal\_date}，必要时加行业等分组 & 不执行排名与标准化 \\
    排序键 & 证券内按观察日升序 & 信号日内按稳定证券标识确定性排序 & 排序必须显式 \\
    典型操作 & \texttt{shift}、\texttt{rolling}、历史端点 & 排名、分位数、截面标准化 & 动量使用证券内历史 \\
    泄漏风险 & 窗口跨证券或包含信号日之后数据 & 混入未来日期或全样本统计量 & 两类错误均应断言失败 \\
    输出检查 & 首尾日期、观测数、窗口完整性 & 每日样本量、并列值和缺失处理 & 保留原值和有效原因 \\
    \bottomrule
  \end{tabularx}
\end{table}

\section{因子定义卡与统一输出接口}

定义卡是研究假设与代码之间的合同。它应在查看后续表现之前冻结，并随版本管理。\cref{tab:ch04-definition-cards} 汇总本章三个教学因子的定义卡；表中的预期方向是待检验假设，不是实证结论。

{
\setlength{\tabcolsep}{3pt}
\footnotesize
\begin{longtable}{p{2.05cm} p{4.1cm} p{4.1cm} p{4.1cm}}
\caption{三个教学因子的综合定义卡}\label{tab:ch04-definition-cards}\\
\toprule
字段 & 中期动量 & 账面市值比价值 & ROE 类质量 \\
\midrule
\endfirsthead
\toprule
字段 & 中期动量 & 账面市值比价值 & ROE 类质量 \\
\midrule
\endhead
因子名称 & \texttt{momentum\_medium} & \texttt{book\_to\_price} & \texttt{roe\_quality\_ttm} \\
研究假设 & 中期价格状态可能具有延续性 & 账面权益相对市值较高可能对应价值暴露 & 盈利相对权益基础较高可能对应质量暴露 \\
经济逻辑 & 用较长历史区间、跳过短期窗口描述中期状态 & 比较点时账面权益与同一信号时点市值 & 比较点时 TTM 利润与平均点时权益 \\
数学公式 & 见 \cref{eq:ch04-momentum} & 见 \cref{eq:ch04-book-to-price} & 见 \cref{eq:ch04-roe,eq:ch04-average-equity} \\
预期方向 & 原始值越高，教学假设方向越正 & 原始值越高，教学假设方向越正 & 原始值越高，教学假设方向越正 \\
输入数据 & 信号截止前价格或总收益序列 & 点时账面权益、信号时点市值 & 点时 TTM 利润、当期与历史权益 \\
可得时间 & 两端点及窗口状态均不晚于 \(t\) & 财务可得日与市值观察日均不晚于 \(t\) & 所有 TTM 组成项和权益记录均不晚于 \(t\) \\
参数 & \(L_{\mathrm{lookback}}\)、\(L_{\mathrm{skip}}\)、最小观测数 & 权益与市值口径、单位规则、分母阈值 & TTM 规则、平均权益规则、分母阈值 \\
最小历史长度 & 覆盖长回看端点并满足最小观测数 & 至少一条有效点时权益与同日市值 & 完整 TTM 组成项及两期可比权益 \\
有效条件 & 同证券、端点存在、历史充足、无未来输入 & 分子分母可得、单位一致、分母合法 & TTM 完整、平均权益合法、历史记录可得 \\
无效原因 & 历史不足、端点缺失、未来输入、公司行动状态待确认 & 财务缺失、非正权益、非法市值、单位不一致 & TTM 不完整、权益缺失、分母非正或过小 \\
输出频率 & 每个周度信号日逐证券一条 & 每个周度信号日逐证券一条 & 每个周度信号日逐证券一条 \\
版本 & 公式、参数与数据口径共同组成版本 & 公式、权益与市值口径共同组成版本 & 公式、TTM 与权益口径共同组成版本 \\
人工核对 & 检查两端点、观测数并手算价格比 & 追溯点时权益和市值后手算比值 & 展开 TTM 组成与平均权益后手算 \\
\bottomrule
\end{longtable}
}

统一输出接口应让三个因子可以纵向拼接，同时保留各自的无效原因。\cref{tab:ch04-output-schema} 是教学基线 schema；实际存储类型由项目实现决定，但字段语义不应随因子变化。

\begin{table}[htbp]
  \centering
  \small
  \caption{统一原始因子输出 schema}
  \label{tab:ch04-output-schema}
  \begin{tabularx}{\textwidth}{p{3.15cm} p{2.25cm} Y Y}
    \toprule
    字段 & 教学类型 & 含义 & 最小约束 \\
    \midrule
    \texttt{security\_id} & 标识 & 稳定证券主键 & 不使用展示代码代替稳定主键 \\
    \texttt{signal\_date} & 日期 & 因子所属信号日 & 必须属于冻结调仓日程 \\
    \texttt{factor\_name} & 字符串 & 定义卡中的固定名称 & 不得因后续表现改名或翻向 \\
    \texttt{raw\_value} & 数值或缺失 & 未经第 5 章处理的原始值 & 无效时保持缺失，不填零 \\
    \texttt{is\_valid} & 布尔 & 是否满足定义卡有效条件 & 与无效原因分开保存 \\
    \texttt{invalid\_reason} & 枚举或列表 & 无法给出原始值的原因 & 有效时为空；优先级需确定 \\
    \texttt{risk\_reason} & 枚举或列表 & 数值可算但需复核的风险 & 可与有效值同时存在 \\
    \texttt{input\_latest\_date} & 日期 & 本次实际使用输入的最晚可得日 & 必须不晚于信号日 \\
    \texttt{factor\_version} & 字符串 & 公式和实现版本 & 版本变更不得覆盖历史输出 \\
    \texttt{config\_version} & 字符串 & 参数与数据口径配置版本 & 可追溯到冻结配置 \\
    \bottomrule
  \end{tabularx}
\end{table}

同一 \texttt{security\_id}、\texttt{signal\_date}、\texttt{factor\_name}、\texttt{factor\_version} 和 \texttt{config\_version} 组合必须唯一。若多项无效条件同时命中，可以保存原因列表，也可以按定义卡规定的确定性优先级选择主原因，但不能由代码执行顺序偶然决定。

\section{中期动量因子的定义与实现}

中期动量用较早长回看端点到较近短期跳过端点之间的价格变化描述证券的中期状态。令 \(t\) 为信号日，\(L_{\mathrm{lookback}}\) 为长回看长度，\(L_{\mathrm{skip}}\) 为短期跳过长度，且
\(L_{\mathrm{lookback}}>L_{\mathrm{skip}}\geq 0\)。教学基线定义为

\begin{equation}
  \mathrm{MOM}_{i,t}
  =
  \frac{
    P^{\mathrm{adj}}_{i,\,t-L_{\mathrm{skip}}}
  }{
    P^{\mathrm{adj}}_{i,\,t-L_{\mathrm{lookback}}}
  }
  - 1 .
  \label{eq:ch04-momentum}
\end{equation}

其中 \(P^{\mathrm{adj}}_{i,s}\) 表示证券 \(i\) 在历史端点 \(s\) 的一致口径价格或总收益代理。两个 \(L\) 应按冻结交易日历或配置中的观察索引解释，不能把自然日与交易日观测混用。实现时先按 \texttt{security\_id} 与观察日排序，再在每只证券内部定位两个端点；任何 \texttt{shift} 或 \texttt{rolling} 都不得跨证券。

\begin{teachingbox}
[教学简化] 候选基线可描述为“约 12 个月回看、跳过最近约 1 个月”。它只是预注册教学设定，不是行业唯一标准，也不是经本项目回测挑出的最优窗口。实际 \(L_{\mathrm{lookback}}\)、\(L_{\mathrm{skip}}\) 和最小观测数应在查看后续标签表现前写入配置。
\end{teachingbox}

只要端点之一缺失、有效历史不足或窗口中出现晚于信号截止的记录，就不应计算原始值。停牌可能导致观测间隔异常，上市历史不足会使长端点不存在，公司行动可能改变价格序列的可比性；这些状态都应进入有效性证据而不是静默前填。

\begin{confirmbox}
[需实际确认] \(P^{\mathrm{adj}}\) 使用何种复权价格或总收益口径、公司行动何时可得、停牌期间如何定义端点、允许的最大日期偏移及最小观测数，都依赖数据源和研究配置。不得凭供应商字段名或最终回测表现决定这些规则。
\end{confirmbox}

\section{账面市值比价值因子的定义与实现}

账面市值比把信号日已经可得的账面权益与同一时点的市值放在同一尺度上比较。令
\(\mathrm{BookEquity}^{\mathrm{PIT}}_{i,t}\) 为按 \cref{ch:data-timing} 生成的点时账面权益快照，\(\mathrm{MarketCap}_{i,t}\) 为与信号时点一致的市值，则

\begin{equation}
  \mathrm{BP}_{i,t}
  =
  \frac{
    \mathrm{BookEquity}^{\mathrm{PIT}}_{i,t}
  }{
    \mathrm{MarketCap}_{i,t}
  } .
  \label{eq:ch04-book-to-price}
\end{equation}

财务分子的报告期可以早于 \(t\)，但其可得日必须不晚于 \(t\)；市值分母也必须对应信号时点，而不是未来日期或另一套未说明的频率。计算前要验证两者币种和数量级一致。若单位不一致，即使除法可以执行，结果也没有可解释性。

本章教学基线把缺失或非正账面权益、缺失或非正市值、极端小分母以及单位不一致列为显式无效状态。这里的“无效”表示不按当前定义卡生成 BP 原始值，不表示该证券从历史股票池中消失，更不允许用零替代未知。

\begin{confirmbox}
[需实际确认] 账面权益采用何种会计口径，市值采用总市值、流通市值或其他定义，币种与单位如何映射，非正权益和小分母阈值如何配置，均由研究目标、数据源和团队口径决定。本章不把任何一种选择写成唯一标准。
\end{confirmbox}

\section{ROE 类质量因子的定义与实现}

ROE 类质量因子用一段期间利润与支持该利润的平均权益比较。令
\(\mathrm{NetIncome}^{\mathrm{TTM,PIT}}_{i,t}\) 表示在信号日已经完整可得的 TTM 利润，令
\(\mathrm{AverageBookEquity}^{\mathrm{PIT}}_{i,t}\) 表示平均点时权益，教学基线为

\begin{equation}
  \mathrm{ROE}^{\mathrm{TTM}}_{i,t}
  =
  \frac{
    \mathrm{NetIncome}^{\mathrm{TTM,PIT}}_{i,t}
  }{
    \mathrm{AverageBookEquity}^{\mathrm{PIT}}_{i,t}
  } ,
  \label{eq:ch04-roe}
\end{equation}

其中平均权益可用一般形式表示为

\begin{equation}
  \mathrm{AverageBookEquity}^{\mathrm{PIT}}_{i,t}
  =
  \frac{
    \mathrm{BookEquity}^{\mathrm{PIT}}_{i,t}
    +
    \mathrm{BookEquity}^{\mathrm{PIT}}_{i,t-1y}
  }{2}.
  \label{eq:ch04-average-equity}
\end{equation}

式中的 \(t-1y\) 表示定义卡要求的可比历史权益期，不意味着从信号日机械减去一个自然年后直接取值。两期权益记录及构成 TTM 利润的每个组成项，都必须在 \(t\) 时点已经可得。不能用后来修订的全年利润替换历史当时可见的组成，也不能在 TTM 项缺失时从未来报告补齐。

负利润可以产生负的 ROE 原始值，并不自动等于数学无效；但平均权益为零、负值或低于配置阈值时，教学基线拒绝计算。若数值可算但一次性项目可能显著影响解释，应保留 \texttt{risk\_reason}，而不是在本章擅自调整利润或缩尾。

\begin{confirmbox}
[需实际确认] 累计年初至今数据如何转换为单季或 TTM、归母利润或其他净利润口径、归母权益或其他权益口径、历史可比期、修订版本、一次性项目识别和小分母阈值，都依赖实际财务数据结构与会计口径。
\end{confirmbox}

\section{三个因子的数据依赖与点时边界}

三个因子使用不同数据，但共享同一条时间约束：实际参与计算的最晚可得日不能晚于信号日 \(t\)。\cref{tab:ch04-data-dependencies} 只列概念字段，不假设任何供应商字段名。

\begin{table}[htbp]
  \centering
  \small
  \caption{三个因子的数据依赖与点时边界}
  \label{tab:ch04-data-dependencies}
  \begin{tabularx}{\textwidth}{p{2.55cm} p{3.4cm} p{3.15cm} Y}
    \toprule
    因子 & 概念输入 & 截止要求 & 关键核对 \\
    \midrule
    中期动量 & 同证券历史价格或总收益序列、交易日历 & 两端点和窗口记录均不晚于 \(t\) & 分组、排序、端点、观测数、公司行动状态 \\
    账面市值比 & 点时账面权益、信号时点市值、单位元数据 & 财务可得日和市值观察日均不晚于 \(t\) & 点时版本、币种、数量级、非正值和小分母 \\
    ROE 类质量 & 点时 TTM 利润组成、当期及历史权益、单位元数据 & 每个利润组成项与两期权益均不晚于 \(t\) & TTM 完整性、历史可比期、修订、平均权益 \\
    \bottomrule
  \end{tabularx}
\end{table}

因子程序的允许输入应来自第 2 章冻结的点时快照与第 3 章冻结的
\(\mathcal{U}^{\mathrm{signal}}_t\)。\texttt{future\_return}、\texttt{execution\_status}、订单成交结果和实际持仓都属于未来或后续层，不得出现在因子计算函数的数据依赖中。即使只是“方便过滤”，访问这些表也会破坏职责边界。

\section{有效值规则、无效原因和缺失状态}

\texttt{is\_valid} 回答“当前定义是否允许使用该数值”，\texttt{invalid\_reason} 回答“为什么不允许”，\texttt{raw\_value} 则只在有效时保存公式输出。三列必须分开：缺失原始值不等于实现错误，布尔值也不能替代原因证据。

建议的教学原因码包括：

\begin{itemize}
  \item 动量：\texttt{MOM\_INSUFFICIENT\_HISTORY}、\texttt{MOM\_ENDPOINT\_MISSING}、\texttt{MOM\_FUTURE\_INPUT}；
  \item 账面市值比：\texttt{BP\_BOOK\_MISSING}、\texttt{BP\_NONPOSITIVE\_EQUITY}、\texttt{BP\_MARKET\_CAP\_INVALID}、\texttt{UNIT\_MISMATCH}；
  \item ROE：\texttt{ROE\_TTM\_INCOMPLETE}、\texttt{ROE\_EQUITY\_MISSING}、\texttt{ROE\_DENOMINATOR\_INVALID}；
  \item 公共检查：\texttt{PIT\_AFTER\_SIGNAL}、\texttt{DUPLICATE\_KEY}、\texttt{CONFIG\_MISSING}。
\end{itemize}

这些机器码是教学字典，不代表团队内部枚举。[机构实现] 最终原因优先级、是否允许多原因、日志位置和输入血缘标识应由项目接口规范确定。无论采用何种形式，都不得用 \(0\) 代替未知、无效或未计算。

\section{一个匿名证券的人工计算案例}

以下虚构数字仅用于解释公式与实现，不代表真实证券、真实数据库或实证研究结果。设唯一信号日为 \(t^\star\)，所有金额均为同一“教学货币单位”；“样本 A”“样本 B”和“样本 C”均为匿名代称。

\begin{table}[htbp]
  \centering
  \small
  \caption{三个匿名证券在单一信号日的教学输入}
  \label{tab:ch04-toy-inputs}
  \begin{tabularx}{\textwidth}{p{1.7cm} p{3.4cm} p{3.2cm} p{3.6cm} Y}
    \toprule
    证券 & 动量端点与历史 & BP 输入 & ROE 输入 & 点时状态 \\
    \midrule
    样本 A & 长端点 100，跳过端点 112；历史充足 & 权益 60，市值 150 & TTM 利润 9；权益 60、40 & 全部不晚于 \(t^\star\) \\
    样本 B & 长端点缺失；历史不足 & 权益 \(-5\)，市值 80 & TTM 利润 2；权益 \(-5\)、5 & 全部不晚于 \(t^\star\) \\
    样本 C & 长端点 100，跳过端点 105 & 权益记录在 \(t^\star\) 后才可得 & TTM 组成项不完整 & 财务输入越过截止日 \\
    \bottomrule
  \end{tabularx}
\end{table}

\begin{casebox}
样本 A 的动量手算为 \(112/100-1=0.12\)；BP 为 \(60/150=0.40\)；平均权益为 \((60+40)/2=50\)，所以 ROE 类原始值为 \(9/50=0.18\)。这些数值只是公式核对结果，不是收益、因子检验或策略表现。

样本 B 因长端点缺失而输出
\texttt{MOM\_INSUFFICIENT\_HISTORY}；其账面权益为负，BP 按教学基线无效；平均权益为零，ROE 输出
\texttt{ROE\_DENOMINATOR\_INVALID}。样本 C 的财务记录在 \(t^\star\) 后才可得，不能用于 BP；TTM 组成项也不得由未来报告补齐。三只证券的信号池记录都应保留，无效因子不等于证券历史不存在。
\end{casebox}

人工核对时应从输出反向追溯到每个输入端点、可得日、单位和配置。程序值与手算值一致只是必要条件；若手算本身使用了未来记录，两者一致仍然是错误实现。

\section{Python/pandas 风格最小实现}

\cref{lst:ch04-factor-interface} 展示接口与防错逻辑。它假设点时快照和信号池已经由前两章冻结，不重新实现数据连接，也不依赖供应商字段。实际项目应把原因码优先级和单位映射放入版本化配置。

\begin{lstlisting}[
  caption={三个原始因子的最小接口与防错逻辑},
  label={lst:ch04-factor-interface}
]
FORBIDDEN = {"future_return", "execution_status", "filled"}
KEY = [
    "security_id", "signal_date", "factor_name",
    "factor_version", "config_version",
]

def check_frozen_inputs(frame):
    assert FORBIDDEN.isdisjoint(frame.columns)
    assert frame["signal_eligible"].all()
    assert frame["pit_snapshot_frozen"].all()
    assert (frame["input_latest_date"]
            <= frame["signal_date"]).all()

def compute_momentum(prices, config):
    assert config["lookback_lag"] > config["skip_lag"]
    assert config["min_observations"] > config["lookback_lag"]
    p = prices.sort_values(["security_id", "observation_date"])
    by_security = p.groupby("security_id", sort=False)
    p["p_skip"] = by_security["price_proxy"].shift(
        config["skip_lag"]
    )
    p["p_lookback"] = by_security["price_proxy"].shift(
        config["lookback_lag"]
    )
    p["history_count"] = by_security.cumcount() + 1
    enough = p["history_count"].ge(config["min_observations"])
    endpoints = p[["p_skip", "p_lookback"]].gt(0).all(axis=1)
    p["is_valid"] = enough & endpoints
    p["raw_value"] = (
        p["p_skip"] / p["p_lookback"] - 1
    ).where(p["is_valid"])
    p["invalid_reason"] = None
    p.loc[~enough, "invalid_reason"] = \
        "MOM_INSUFFICIENT_HISTORY"
    p.loc[enough & ~endpoints, "invalid_reason"] = \
        "MOM_ENDPOINT_MISSING"
    return p

def compute_book_to_price(snapshot, config):
    out = snapshot.copy()
    book_present = snapshot["book_equity_pit"].notna()
    market_present = snapshot["market_cap"].notna()
    unit_ok = snapshot["book_unit"].eq(snapshot["market_cap_unit"])
    denom_ok = market_present & snapshot["market_cap"].gt(
        config["denom_floor"]
    )
    equity_ok = book_present & snapshot["book_equity_pit"].gt(0)
    valid = book_present & unit_ok & denom_ok & equity_ok
    out["raw_value"] = (
        snapshot["book_equity_pit"] / snapshot["market_cap"]
    ).where(valid)
    out["is_valid"] = valid
    out["invalid_reason"] = None
    out.loc[~book_present, "invalid_reason"] = "BP_BOOK_MISSING"
    out.loc[book_present & ~unit_ok, "invalid_reason"] = \
        "UNIT_MISMATCH"
    out.loc[book_present & unit_ok & ~denom_ok, "invalid_reason"] = \
        "BP_MARKET_CAP_INVALID"
    out.loc[book_present & unit_ok & denom_ok & ~equity_ok,
            "invalid_reason"] = \
        "BP_NONPOSITIVE_EQUITY"
    return out

def compute_roe_quality(snapshot, config):
    out = snapshot.copy()
    equity_complete = snapshot[
        ["book_equity_pit", "book_equity_prior_pit"]
    ].notna().all(axis=1)
    average_equity = (
        snapshot["book_equity_pit"]
        + snapshot["book_equity_prior_pit"]
    ) / 2
    ttm_ok = (
        snapshot["ttm_complete"]
        & snapshot["net_income_ttm_pit"].notna()
    )
    denom_ok = equity_complete & average_equity.gt(
        config["equity_floor"]
    )
    valid = ttm_ok & equity_complete & denom_ok
    out["raw_value"] = (
        snapshot["net_income_ttm_pit"] / average_equity
    ).where(valid)
    out["is_valid"] = valid
    out["invalid_reason"] = None
    out.loc[~ttm_ok, "invalid_reason"] = "ROE_TTM_INCOMPLETE"
    out.loc[ttm_ok & ~equity_complete, "invalid_reason"] = \
        "ROE_EQUITY_MISSING"
    out.loc[ttm_ok & equity_complete & ~denom_ok, "invalid_reason"] = \
        "ROE_DENOMINATOR_INVALID"
    return out

def finalize_output(output):
    check_frozen_inputs(output)
    required = set(KEY + [
        "raw_value", "is_valid", "invalid_reason",
        "input_latest_date",
    ])
    assert required.issubset(output.columns)
    assert not output.duplicated(KEY).any()
    assert output.loc[~output["is_valid"], "raw_value"].isna().all()
    return output  # no fillna(0), winsorize, rank, or neutralize
\end{lstlisting}

代码中两次 \texttt{shift} 都来自同一个
\texttt{groupby("security\_id")} 对象，因此不会跨证券。示例没有展开所有原因码分支，但生产实现必须把每个失败规则映射到
\texttt{invalid\_reason}。统一入口还应先把冻结信号池与因子输入按稳定主键连接，且只能对
\(\mathcal{U}^{\mathrm{signal}}_t\) 中的记录生成输出。

\section{单元测试、边界测试和人工穿行核对}

测试的目的不是证明因子有预测力，而是让错误实现尽早失败。自动化测试至少应覆盖：

\begin{enumerate}
  \item 构造两只证券相邻记录，证明未按 \texttt{security\_id} 分组的动量 \texttt{shift} 会被测试捕获；
  \item 窗口内任一 \texttt{input\_latest\_date} 晚于信号日时触发失败；
  \item 历史观测不足时输出 \texttt{MOM\_INSUFFICIENT\_HISTORY}，而不是数值或零；
  \item 财务可得日晚于信号日时，不得参与 BP 或 ROE 计算；
  \item 市值与权益单位不一致时触发 \texttt{UNIT\_MISMATCH}；
  \item 账面权益为负、零或缺失时，按冻结配置输出明确状态；
  \item ROE 平均权益为零、负值或低于阈值时输出明确原因；
  \item TTM 组成项缺失时，不得使用未来报告补齐；
  \item 相同主键与两个版本字段重复时触发唯一性错误；
  \item 因子名称和预期方向与定义卡不一致时失败，不得按后续表现翻转；
  \item 匿名案例手算值与程序值在约定数值精度内一致；
  \item 静态或接口测试证明因子代码不能访问标签表与成交状态表。
\end{enumerate}

\begin{practicebox}
人工穿行至少选择五类记录：一只三个因子均正常的证券；一只上市历史不足的证券；一只财务记录在信号日尚不可得的证券；一只负权益或小分母证券；以及一只发生代码变化或输入缺失的证券。逐条记录信号日、稳定主键、输入观察日与可得日、公式参数、手算结果、程序结果、有效状态、原因码、配置版本和核验人备注。
\end{practicebox}

对动量还应人工查看证券切换边界前后的若干行，确认窗口没有串号；对财务因子应展开每条点时记录的报告期、可得日与修订版本。若测试只验证最终数值而不检查输入证据，仍可能遗漏前视偏差。

\section{常见错误}

\begin{warningbox}
常见错误包括：未排序就计算 \texttt{shift}；滚动窗口跨证券；把自然日长度当作固定交易观测数；用信号日之后价格补齐动量端点；把报告期末当作财务可得日；混用权益和市值单位；用未来全年值拼 TTM；分母为零仍执行除法；将缺失值填零；在原始构造阶段提前缩尾或排名；按后续结果翻转方向或挑选窗口；以及让因子函数读取未来收益、成交状态或实际持仓。
\end{warningbox}

这些错误有的会直接报错，有的却会生成平滑且“看起来合理”的数值。后者更危险，因为它们可能系统性改变覆盖率和横截面排序。修复时应从定义卡、输入血缘和最小匿名案例开始，而不是通过观察最终回测结果猜测原因。

\section{本章产出}

\begin{outputbox}
本章应提交：三张冻结因子定义卡；三个原始因子的公式与参数配置；统一输出 schema 和原因码字典；逐信号日原始因子快照；最小实现接口；自动化单元与边界测试结果；以及覆盖正常、历史不足、点时不可得、非法分母和代码变化或缺失输入的人工穿行记录。
\end{outputbox}

这些产出不包含预处理值、因子评价、参数择优或组合结果。仓库当前没有真实金融数据，因此本章不生成任何实证数值、图形或表现结论。

\section{章节自测}

\begin{quizbox}
\begin{enumerate}
  \item 为什么 \texttt{raw\_value} 不能被排名值或标准化值覆盖？
  \item 动量的两个端点分别由哪个参数决定？为什么必须按证券分组？
  \item “约 12 个月回看、跳过约 1 个月”属于通用原则还是教学基线？
  \item BP 的财务分子和市值分母分别需要满足什么时间约束？
  \item 为什么单位不一致比单纯缺失更难从最终数值中发现？
  \item TTM 利润缺少一个组成项时，为什么不能等待未来报告再回填？
  \item 负利润与负平均权益对 ROE 有效性的含义是否相同？
  \item \texttt{is\_valid}、\texttt{invalid\_reason} 和
  \texttt{risk\_reason} 各自回答什么问题？
  \item 如何证明因子程序没有读取未来收益或成交状态？
  \item 哪五个字段共同构成同一因子输出的教学唯一键？
\end{enumerate}
\end{quizbox}

\section{与第 5 章《因子预处理与基础暴露控制》的衔接}

本章输出的是带有效状态和完整血缘的原始因子长表。第 5 章只能在
\texttt{is\_valid=True} 且口径已经确认的原始值上讨论缺失机制、异常值、变换、排名、标准化以及行业或市值暴露控制；它不应回头改变本章公式、点时输入或因子方向。

衔接时应保留原始列和版本，另建处理后字段，并记录每一步参数与前后样本数。这样，后续任何分布变化都能区分为原始数据变化、有效性规则变化或预处理变化，而不会把三种原因混在一起。

\chapter{因子预处理与基础暴露控制}

\section*{待编写}

% TODO: 按 P03 后续里程碑要求编写本章。

\chapter{单因子检验与可交易性诊断}

\section*{待编写}

% TODO: 按 P03 后续里程碑要求编写本章。

\chapter{稳健性、样本外与反过拟合}

\section*{待编写}

% TODO: 按 P03 后续里程碑要求编写本章。

\chapter{因子冗余与简单多因子合成}

\section*{待编写}

% TODO: 按 P03 后续里程碑要求编写本章。

\chapter{从综合得分到约束型多头组合}
\label{ch:portfolio}

\section{本章导读与学习目标}

第 \cref{ch:factor-combination} 已经输出唯一冻结的 \texttt{combination\_version}、证券层面综合得分、共同样本状态和适用条件。本章把这份横截面排序信息转化为信号日结束后“希望持有”的多头目标权重，并生成相对上一期实际持仓的期望交易差额。核心问题是：\keyterm{如何在不偷看下一交易日的前提下，让目标组合既表达综合分，又满足基础约束并可被完整审计？}

完成本章后，读者应能够：

\begin{enumerate}
  \item 区分综合得分、目标组合、期望订单、实际成交和实际持仓；
  \item 用冻结的固定数量规则和确定性并列规则构造候选集合；
  \item 从基础等权出发，依次处理个股、行业、流动性和换手约束；
  \item 解释持仓缓冲、现金权重与目标换手的作用和边界；
  \item 在约束冲突时执行预注册、确定性的降级流程；
  \item 输出可追溯的目标权重、期望交易差额、约束状态和组合版本；
  \item 为第 10 章保留目标与实际之间的差异，而不提前模拟成交。
\end{enumerate}

\begin{principlebox}
\textbf{[通用原则]} 下一交易日无法成交不能反向改变信号日综合分、候选名单或原始目标权重。未成交和实际持仓偏离应由第 10 章的订单、成交与持仓账本处理。
\end{principlebox}

\begin{figure}[htbp]
  \centering
  \begin{tikzpicture}[
    node distance=2.5mm,
    every node/.style={font=\footnotesize},
    quantPortfolioScore/.style={quantBlueCard, text width=2.15cm, minimum height=1.15cm, align=center},
    quantPortfolioCandidate/.style={quantCard, text width=2.15cm, minimum height=1.15cm, align=center},
    quantPortfolioBase/.style={quantCard, text width=2.15cm, minimum height=1.15cm, align=center},
    quantPortfolioConstraint/.style={quantOrangeCard, text width=2.15cm, minimum height=1.15cm, align=center},
    quantPortfolioTarget/.style={quantGreenCard, text width=2.15cm, minimum height=1.15cm, align=center},
    quantPortfolioTrade/.style={quantBlueCard, text width=2.15cm, minimum height=1.15cm, align=center}
  ]
    \node[quantPortfolioScore] (score) {冻结综合分};
    \node[quantPortfolioCandidate, right=of score] (candidate) {候选集合};
    \node[quantPortfolioBase, right=of candidate] (base) {基础等权};
    \node[quantPortfolioConstraint, right=of base] (constraint) {个股、行业、流动性、换手};
    \node[quantPortfolioTarget, right=of constraint] (target) {目标权重与现金};
    \node[quantPortfolioTrade, right=of target] (trade) {期望交易差额};

    \draw[quantArrow] (score) -- (candidate);
    \draw[quantArrow] (candidate) -- (base);
    \draw[quantArrow] (base) -- (constraint);
    \draw[quantArrow] (constraint) -- (target);
    \draw[quantArrow] (target) -- (trade);
  \end{tikzpicture}
  \caption{从冻结综合分到期望交易差额的规则型流程（结构示意）}
  \label{fig:ch09-portfolio-flow}
\end{figure}

本章不判断下一交易日真实成交，不处理停牌、价格限制、部分成交、正式交易成本、实际持仓更新或组合收益。这些全部留给第 10 章。

\section{为什么综合得分不是持仓}

\texttt{combined\_score} 只表达证券在冻结因子体系下的相对优先级，没有资金单位，也没有自动包含行业、流动性、换手和个股集中度。把分数直接归一化为权重，可能产生负权重、集中于少数证券，或让某一共同暴露被放大。

目标权重是一项带约束的决策。它既要尽可能保留排序信息，又要接受可行性边界。约束可能削弱原始信号表达；这种损失应通过基础权重、每一步调整和最终状态被记录，而不是隐藏在一个最终权重字段中。

\section{综合分、目标权重、订单和实际持仓的区别}

\begin{table}[htbp]
  \centering
  \small
  \caption{五类对象的职责与时间边界}
  \label{tab:ch09-object-responsibilities}
  \begin{tabularx}{\textwidth}{p{2.7cm} p{4.5cm} p{3.2cm} Y}
    \toprule
    对象 & 表达含义 & 形成时点 & 本章是否生成 \\
    \midrule
    综合得分 & 证券横截面排序信息 & 信号日数据冻结后 & 读取，不改写 \\
    目标组合 & 信号日后希望持有的权重 & 约束构建完成后 & 是 \\
    期望订单 & 目标与上一期实际持仓之差 & 目标权重确定后 & 是，仅目标差额 \\
    实际成交 & 下一交易日真实执行结果 & 交易状态与价格可得后 & 否 \\
    实际持仓 & 成交后真实持有状态 & 成交账本更新后 & 否 \\
    \bottomrule
  \end{tabularx}
\end{table}

期望交易为正表示希望增加，为负表示希望减少，为零表示目标不变，但都不代表已成交。若卖单在下一交易日无法成交，旧持仓仍保留在实际账本中；不能因目标权重为零就让它凭空消失。

\section{本章冻结输入和禁止访问的数据}

本章只读取第 8 章冻结的 \texttt{combination\_version} 及其 \texttt{signal\_date}、\texttt{security\_id}、\texttt{combined\_score}、\texttt{joint\_valid} 和 \texttt{applicability\_conditions}；同时读取 \cref{ch:universe-labels} 的信号日股票池、信号日点时行业、市值与历史流动性，以及上一期实际持仓或期初快照。

禁止访问的字段包括 \texttt{execution\_status}、\texttt{filled}、\texttt{execution\_price}、\texttt{limit\_status}、\texttt{next\_day\_suspension}、\texttt{future\_return} 和 \texttt{backtest\_result}。若输入表意外携带这些字段，目标构建应立即失败，而不是简单忽略后继续运行。

\begin{confirmbox}
\textbf{[需实际确认]} 真实资金规模、基准行业权重、行业偏离、参与率、个股上限、持股数量、换手预算和持仓快照截止时点必须由版本化配置确认。本章不虚构团队阈值。
\end{confirmbox}

\section{从综合分到候选集合}

在信号日 \(t\)，先保留 \texttt{joint\_valid = True} 且满足冻结适用条件的证券，再按 \texttt{combined\_score} 从高到低排序。教学主基线使用预注册的固定数量 \(N\)，候选集合可写为

\begin{equation}
  C_t
  = \operatorname{TopN}_{i\in U_t^{\mathrm{eligible}}}
    \left(\operatorname{combined\_score}_{i,t};N\right).
  \label{eq:ch09-candidate-set}
\end{equation}

并列时依次使用冻结得分、预注册的次级排序字段和稳定证券主键。稳定主键只用于确定性打破并列，不表达经济偏好。每条记录都保存 \texttt{score\_rank}、\texttt{selected} 与 \texttt{selection\_reason}，未入选记录仍保留。

\section{固定数量与固定分位入选规则}

\begin{table}[htbp]
  \centering
  \small
  \caption{固定数量与固定分位规则的比较}
  \label{tab:ch09-selection-rules}
  \begin{tabularx}{\textwidth}{p{2.6cm} Y Y}
    \toprule
    比较项 & 固定数量 & 固定分位 \\
    \midrule
    持股数量 & 较稳定，便于解释与审计 & 随有效股票池规模变化 \\
    信号池变化 & 入选比例随池规模变化 & 入选比例相对稳定 \\
    小样本风险 & 候选不足时直接返回状态 & 分位边界和取整更敏感 \\
    并列处理 & 第 \(N\) 名附近需确定性规则 & 分位阈值附近需确定性规则 \\
    本章角色 & 教学主基线 & 预注册敏感性方案 \\
    \bottomrule
  \end{tabularx}
\end{table}

持股数量必须在查看第 10 章正式回测前冻结。不能依据净收益、成本后表现或成交成功率在两种规则之间择优。

\section{基础等权组合}

设候选集合为 \(C_t\)，约束前基础权重为

\begin{equation}
  w^{\mathrm{base}}_{i,t}
  =
  \begin{cases}
    1/|C_t|, & i\in C_t,\\
    0, & i\notin C_t.
  \end{cases}
  \label{eq:ch09-base-weight}
\end{equation}

基础等权只是可审计起点，不是最终目标。即使某只证券得分特别高，也不能临时提高其基础权重。本章既不重新优化第 8 章因子权重，也不覆盖 \texttt{combined\_score}；基础和最终权重必须分列保存。

\section{个股最大权重}

多头、不加杠杆和个股上限可写为

\begin{align}
  w_{i,t} &\ge 0, \\
  \sum_i w_{i,t} &\le 1, \\
  w_{i,t} &\le w_{\max}.
  \label{eq:ch09-long-cap}
\end{align}

\(w_{\max}\) 来自冻结配置。若基础等权已经超过上限，说明候选数量和约束不能同时充分投资。截断超限权重后，任何重新分配都必须再次检查上限；不能先截断，再无条件归一化到 1，导致原本合格的证券重新超限。

\section{行业偏离约束}

令 \(g(i,t)\) 为信号日已生效的历史行业，组合行业权重为

\begin{equation}
  W_{g,t}=\sum_{i:g(i,t)=g}w_{i,t}.
  \label{eq:ch09-industry-weight}
\end{equation}

若存在信号日可得的基准行业权重 \(b_{g,t}\)，则可要求

\begin{equation}
  \left|W_{g,t}-b_{g,t}\right|\le \delta_g.
  \label{eq:ch09-industry-deviation}
\end{equation}

历史行业和基准版本都不得晚于信号日。若没有可靠历史基准权重，可采用预注册行业配额或行业内选股作为教学替代，但必须使用独立约束版本。行业约束可能削弱高分证券的表达，这是需要报告的约束代价。

\section{流动性约束}

流动性只能使用信号日前已可得的历史代理。一般目标名义金额上限可写为

\begin{equation}
  \operatorname{TargetNotional}_{i,t}
  \le
  \operatorname{ParticipationCap}
  \times
  \operatorname{HistoricalLiquidity}_{i,t}.
  \label{eq:ch09-liquidity-cap}
\end{equation}

历史窗口必须截止于信号日。资金规模、参与率、价格与成交额口径需实际确认。本章只把代理转换为目标权重上限；下一交易日是否真的成交由第 10 章处理，不能使用下一交易日成交量或不可买状态反向限制目标。

\section{换手预算}

设上一期实际证券权重为 \(h^-_{i,t}\)，目标权重为 \(w_{i,t}\)。教学口径的单边目标换手定义为

\begin{equation}
  \operatorname{Turnover}^{\mathrm{target}}_t
  = \frac{1}{2}\sum_i\left|w_{i,t}-h^-_{i,t}\right|.
  \label{eq:ch09-target-turnover}
\end{equation}

上一期实际持仓在信号日是已知状态，可以作为输入；目标换手不是实际成交换手。若超过冻结预算 \(B_t\)，可按固定比例收缩目标变化或保留部分旧持仓，随后重新检查个股、行业和流动性硬约束。下一交易日未成交信息不能改变本章目标换手。

\section{持仓缓冲与进入/退出阈值}

缓冲区让新证券只有进入更高阈值才买入，让已持有证券只有跌出较低阈值才退出；进入阈值高于退出阈值，可减少边界附近反复交易。缓冲会让目标组合依赖上一期实际持仓，但这不是未来信息。

缓冲参数必须在回测前冻结，并比较加入前后的信号表达、目标换手和集中度。本章将其作为可选教学方案，不作为主基线。不能根据第 10 章的换手和净收益反复调整阈值。

\section{现金权重}

证券目标权重之和不必在不可行时强制等于 1。现金权重定义为

\begin{equation}
  w^{\mathrm{cash}}_t=1-\sum_i w_{i,t},\qquad w^{\mathrm{cash}}_t\ge 0.
  \label{eq:ch09-cash-weight}
\end{equation}

现金是约束后的合法结果，不是自动报错。若个股、行业、流动性或换手限制阻止充分投资，应增加现金并记录 \texttt{PARTIAL\_ALLOCATION}，而不是放宽硬约束或强行归一化。

\section{多项约束的优先级}

\begin{table}[htbp]
  \centering
  \small
  \caption{规则型组合构建的约束优先级}
  \label{tab:ch09-priority}
  \begin{tabularx}{\textwidth}{p{1.4cm} p{4.6cm} Y}
    \toprule
    优先级 & 约束或目标 & 违反时处理 \\
    \midrule
    1 & 数据、主键和版本合法 & 立即失败，不生成目标 \\
    2 & 不做空、不加杠杆 & 硬约束，不放宽 \\
    3 & 个股权重上限 & 缩减并再次验证 \\
    4 & 已确认行业约束 & 按固定顺序补充或缩减 \\
    5 & 已确认流动性上限 & 降低目标，不使用未来成交数据 \\
    6 & 换手预算 & 比例收缩或保留旧持仓 \\
    7 & 尽可能表达综合分排序 & 在硬约束可行域内执行 \\
    8 & 充分投资 & 不能压倒更高优先级约束 \\
    9 & 剩余部分持有现金 & 记录部分配置状态 \\
    \bottomrule
  \end{tabularx}
\end{table}

同时适用的基础约束汇总见 \cref{tab:ch09-constraint-summary}。

\begin{table}[htbp]
  \centering
  \footnotesize
  \caption{个股、行业、流动性和换手约束摘要}
  \label{tab:ch09-constraint-summary}
  \begin{tabularx}{\textwidth}{p{2.5cm} p{4.0cm} p{3.6cm} Y}
    \toprule
    约束 & 信号日输入 & 目标检查 & 禁止使用 \\
    \midrule
    个股 & 冻结 \(w_{\max}\) & 每只证券目标权重 & 历史收益择优上限 \\
    行业 & 点时行业、点时基准或配额 & 行业权重与偏离 & 未来行业或基准版本 \\
    流动性 & 信号日前历史代理与资金规模 & 目标名义金额上限 & 下一日成交量和成交状态 \\
    换手 & 上一期实际持仓、冻结预算 & 目标差额口径 & 实际成交或未成交信息 \\
    \bottomrule
  \end{tabularx}
\end{table}

\section{约束冲突、不可行与确定性降级}

约束冲突不能靠随机删股或静默放宽解决。所有降级动作都要按配置顺序执行，并记录输入状态、触发原因、调整前后权重和仍未满足的约束。

\begin{figure}[htbp]
  \centering
  \begin{tikzpicture}[
    node distance=4mm,
    every node/.style={font=\small},
    quantPortfolioDecision/.style={diamond, aspect=2.4, draw=WarmOrange, fill=WarmOrange!8, align=center, inner sep=2pt},
    quantPortfolioAction/.style={quantCard, text width=7.2cm, minimum height=0.9cm, align=center},
    quantPortfolioSuccess/.style={quantGreenCard, text width=7.2cm, minimum height=0.9cm, align=center},
    quantPortfolioFailure/.style={quantOrangeCard, text width=7.2cm, minimum height=0.9cm, align=center}
  ]
    \node[quantPortfolioDecision] (feasible) {全部硬约束可行？};
    \node[quantPortfolioSuccess, below left=7mm and 20mm of feasible] (success) {是：冻结目标权重与约束报告};
    \node[quantPortfolioAction, below right=7mm and 20mm of feasible] (reduce) {否：按高分顺序补充或缩减超限权重};
    \node[quantPortfolioAction, below=of reduce] (turnover) {按冻结规则收缩换手并保留必要旧持仓};
    \node[quantPortfolioAction, below=of turnover] (cash) {无法充分分配的部分转为现金};
    \node[quantPortfolioDecision, below=of cash] (recheck) {重新验证硬约束？};
    \node[quantPortfolioSuccess, below left=7mm and 20mm of recheck] (partial) {通过：\texttt{PARTIAL\_ALLOCATION}};
    \node[quantPortfolioFailure, below right=7mm and 20mm of recheck] (failed) {仍失败：\texttt{INFEASIBLE}};

    \draw[quantArrow] (feasible) -- node[above left]{是} (success);
    \draw[quantArrow] (feasible) -- node[above right]{否} (reduce);
    \draw[quantArrow] (reduce) -- (turnover);
    \draw[quantArrow] (turnover) -- (cash);
    \draw[quantArrow] (cash) -- (recheck);
    \draw[quantArrow] (recheck) -- node[above left]{是} (partial);
    \draw[quantArrow] (recheck) -- node[above right]{否} (failed);
  \end{tikzpicture}
  \caption{约束不可行时的确定性降级决策树（结构示意）}
  \label{fig:ch09-degradation-tree}
\end{figure}

允许的动作包括从下一名高分证券补充、缩减超限权重、按固定顺序重新分配、比例收缩交易、保留部分旧持仓和增加现金。任何动作后都必须重复验证高优先级硬约束。

\section{目标权重与期望交易量}

目标交易差额定义为

\begin{equation}
  \operatorname{desired\_trade\_weight}_{i,t}
  = w_{i,t}-h^-_{i,t}.
  \label{eq:ch09-desired-trade}
\end{equation}

正值表示期望增加，负值表示期望减少，零表示目标不变。若需要目标名义金额，可在冻结资金规模下转换，但仍不能称为已下单或已成交。本章不生成 \texttt{filled\_weight} 或 \texttt{actual\_holding\_after\_trade}。

\section{组合约束检查与审计表}

\begin{table}[htbp]
  \centering
  \small
  \caption{目标组合约束检查模板}
  \label{tab:ch09-constraint-audit}
  \begin{tabularx}{\textwidth}{p{3.2cm} p{3.1cm} p{3.1cm} Y}
    \toprule
    检查项 & 冻结限制 & 目标观察值 & 状态与降级记录 \\
    \midrule
    证券权重非负 & \(w_{i,t}\ge0\) & -- & -- \\
    证券权重与现金 & 合计为 1 & -- & -- \\
    最大个股权重 & \(w_{\max}\) & -- & -- \\
    行业偏离 & \(\delta_g\) 或配额 & -- & -- \\
    流动性上限 & 历史代理映射上限 & -- & -- \\
    目标换手 & \(B_t\) & -- & -- \\
    证券数量与覆盖 & 配置与状态规则 & -- & -- \\
    组合版本 & 唯一且不可覆盖 & -- & -- \\
    \bottomrule
  \end{tabularx}
\end{table}

审计报告还应保存各行业权重及偏离、流动性覆盖、约束失败、每步重新分配、现金来源和最终 \texttt{portfolio\_status}。只保留最终权重会丢失约束如何影响信号的证据链。

\section{一个匿名教学案例}

\begin{casebox}
\textbf{仅用于解释选股、权重和约束逻辑，不代表真实证券、真实数据库、真实投资组合或实证研究结果。}

匿名信号日 \(t^\star\) 有九只证券，教学配置固定选择前四名。上一期证券权重合计为 0.80，现金为 0.20。案例使用匿名约束：个股上限、行业配额和流动性目标上限只为解释计算，不代表真实团队阈值。

\begin{center}
\footnotesize
\begin{tabular}{c c r r c r}
  \toprule
  证券 & 行业 & 综合分 & 上期实际权重 & 前四名 & 约束提示 \\
  \midrule
  A & 甲 & 2.1 & 0.15 & 是 & 高分，受个股上限 \\
  B & 甲 & 1.8 & 0.10 & 是 & 行业甲接近上限 \\
  C & 乙 & 1.6 & 0.08 & 是 & 流动性目标上限较低 \\
  D & 丙 & 1.2 & 0.12 & 是 & 个股上限 \\
  E & 乙 & 0.9 & 0.10 & 否 & 退出目标 \\
  F & 丙 & 0.7 & 0.08 & 否 & 退出目标 \\
  G & 甲 & 0.4 & 0.07 & 否 & 退出目标 \\
  H & 乙 & 0.1 & 0.05 & 否 & 退出目标 \\
  I & 丙 & -0.2 & 0.05 & 否 & 退出目标 \\
  \bottomrule
\end{tabular}
\end{center}

前四名的基础等权均为 0.25。确定性约束后，A 因个股上限降至 0.20；行业甲总额限制使 B 只能为 0.10；C 的历史流动性映射上限为 0.10；D 受个股上限约束为 0.20。证券目标权重合计 0.60，现金为 0.40。

\begin{center}
\footnotesize
\begin{tabular}{c r r r r}
  \toprule
  证券 & 基础权重 & 上期实际权重 & 目标权重 & 期望交易差额 \\
  \midrule
  A & 0.25 & 0.15 & 0.20 & 0.05 \\
  B & 0.25 & 0.10 & 0.10 & 0.00 \\
  C & 0.25 & 0.08 & 0.10 & 0.02 \\
  D & 0.25 & 0.12 & 0.20 & 0.08 \\
  E & 0.00 & 0.10 & 0.00 & -0.10 \\
  F & 0.00 & 0.08 & 0.00 & -0.08 \\
  G & 0.00 & 0.07 & 0.00 & -0.07 \\
  H & 0.00 & 0.05 & 0.00 & -0.05 \\
  I & 0.00 & 0.05 & 0.00 & -0.05 \\
  \bottomrule
\end{tabular}
\end{center}
\end{casebox}

案例的目标单边换手按 \cref{eq:ch09-target-turnover} 为 0.25。表中目标权重和期望差额均不代表下一交易日实际结果；若某笔交易无法成交，第 10 章会保留旧实际持仓并记录偏离，而不是回写本表。

\section{pandas 风格最小实现}

\cref{lst:ch09-portfolio-builder} 展示规则型目标组合构建的职责边界。代码只生成目标和审计状态，不生成订单成交或回测结果。

\begin{lstlisting}[
  caption={约束型多头目标组合的最小接口},
  label={lst:ch09-portfolio-builder}
]
FORBIDDEN_COLUMNS = {
    "execution_status", "filled", "execution_price",
    "limit_status", "next_day_suspension",
    "future_return", "backtest_result", "orders", "fills",
}

def build_constrained_long_only_target(
    frozen_scores,
    signal_universe,
    previous_actual_holdings,
    portfolio_config,
):
    assert portfolio_config.combination_version_frozen
    assert (
        frozen_scores["combination_version"].unique().tolist()
        == [portfolio_config.combination_version]
    )
    for frame in (
        frozen_scores,
        signal_universe,
        previous_actual_holdings,
    ):
        assert FORBIDDEN_COLUMNS.isdisjoint(frame.columns)

    assert_unique(frozen_scores, [
        "security_id", "signal_date", "combination_version"
    ])
    validate_point_in_time_columns(
        signal_universe,
        asof_columns=[
            "industry_asof", "benchmark_weight_asof",
            "liquidity_window_end",
        ],
        not_after="signal_date",
    )

    conditions_met = evaluate_frozen_applicability_conditions(
        frozen_scores["applicability_conditions"],
        signal_universe,
        asof="signal_date",
    )
    eligible = frozen_scores.loc[
        frozen_scores["joint_valid"] & conditions_met
    ].copy()
    ranked = deterministic_rank(
        eligible,
        by=["combined_score", "security_id"],
        ascending=[False, True],
        tie_rule=portfolio_config.tie_rule,
    )
    selected = select_from_frozen_rule(
        ranked,
        method=portfolio_config.selection_method,
        count=portfolio_config.fixed_count,
    )
    if selected.empty:
        return append_infeasible_portfolio(
            reason="NO_ELIGIBLE_CANDIDATE",
            portfolio_version=hash_frozen_config(portfolio_config),
        )
    ranked["selected"] = ranked["security_id"].isin(
        selected["security_id"]
    )
    ranked["base_weight"] = 0.0
    ranked.loc[ranked["selected"], "base_weight"] = (
        1.0 / len(selected)
    )

    state = build_portfolio_state_union(
        ranked,
        signal_universe,
        previous_actual_holdings,
        stable_key="security_id",
        include_previous_only_holdings=True,
    )
    state = apply_optional_buffer(
        state,
        enabled=portfolio_config.buffer_enabled,
        entry_rule=portfolio_config.entry_rule,
        exit_rule=portfolio_config.exit_rule,
    )
    state = apply_individual_caps(
        state, cap=portfolio_config.individual_cap
    )
    state = apply_industry_constraints(
        state,
        limits=portfolio_config.industry_limits,
        deterministic_order=ranked["security_id"],
    )
    state = apply_liquidity_caps(
        state,
        capital=portfolio_config.capital,
        participation_cap=portfolio_config.participation_cap,
    )
    state = redistribute_with_frozen_order(
        state,
        ranked=ranked,
        allow_cash=True,
    )
    validate_hard_constraints(state)

    state = shrink_to_turnover_budget(
        state,
        previous_weight="previous_actual_weight",
        budget=portfolio_config.turnover_budget,
        trace=True,
    )
    state = deterministic_degradation(
        state,
        priority=portfolio_config.constraint_priority,
        allow_cash=True,
        statuses=["PARTIAL_ALLOCATION", "INFEASIBLE"],
    )
    validate_hard_constraints(state)

    state["target_weight"] = state["constrained_weight"]
    state["desired_trade_weight"] = (
        state["target_weight"]
        - state["previous_actual_weight"]
    )
    cash_weight = 1.0 - state["target_weight"].sum()
    assert (state["target_weight"] >= 0).all()
    assert cash_weight >= 0
    assert abs(state["target_weight"].sum() + cash_weight - 1) < 1e-12
    assert FORBIDDEN_COLUMNS.isdisjoint(state.columns)

    output = attach_portfolio_audit_fields(
        state,
        cash_weight=cash_weight,
        portfolio_version=hash_frozen_config(portfolio_config),
    )
    portfolio_store.append_only(output)
    return output
\end{lstlisting}

真实实现还应保证候选不足时不发生除零，并让所有降级步骤追加到审计表。复杂二次规划不是本章要求；规则型流程的价值在于每次调整都能重放和解释。

\section{自动测试与人工核对}

自动测试至少覆盖：

\begin{enumerate}
  \item 非冻结 \texttt{combination\_version} 不能使用；
  \item \texttt{joint\_valid = False} 的记录不能进入候选；
  \item 相同证券、信号日和版本重复时失败；
  \item 排序不能读取未来标签；
  \item 并列处理确定且可重复；
  \item 固定数量或固定分位规则严格来自配置；
  \item 任一证券目标权重不得为负；
  \item 证券目标权重与现金权重合计为 1；
  \item 任一个股不得超过上限；
  \item 重新分配后再次检查个股上限；
  \item 历史行业区间覆盖信号日；
  \item 基准行业权重版本不晚于信号日；
  \item 流动性窗口不包含未来数据；
  \item 下一交易日成交量或状态不得参与目标构建；
  \item 目标换手与实际成交换手使用不同字段名；
  \item 换手预算降级过程可追溯；
  \item 无法充分投资时现金权重为正；
  \item 不可行状态不能被静默忽略；
  \item \texttt{target\_weight}、\texttt{desired\_trade\_weight} 与 \texttt{previous\_actual\_weight} 可相互还原；
  \item 组合版本变化保留旧输出；
  \item 相同输入、配置和版本得到相同结果；
  \item 目标组合函数不能访问 \texttt{orders}、\texttt{fills}、\texttt{execution\_price} 或 \texttt{backtest\_result}。
\end{enumerate}

\begin{practicebox}
\textbf{[正确实践]} 人工至少核对一个普通候选、一个并列边界、一个个股超限、一个行业受限、一个流动性受限、一个换手收缩和一个现金配置案例。逐项重算基础权重、每步约束、目标权重、期望差额和组合状态，并确认没有读取下一交易日字段。
\end{practicebox}

\section{常见错误}

\begin{warningbox}
常见错误包括：把综合分直接当权重；根据第 10 章收益选择持股数量；使用下一交易日不可买状态重排信号日候选；用未来基准权重或行业分类；先截断个股权重再归一化导致重新超限；行业约束后不复查个股约束；为满仓违反流动性限制；把目标换手称为实际换手；把目标组合称为实际持仓；无法卖出时让旧持仓凭空消失；不可行时静默放宽约束；只保存最终权重而不保存基础权重和降级过程；把现金视为回测错误；以及使用组合结果回头修改第 8 章因子集合或权重。
\end{warningbox}

这些错误会把信号、决策和执行三层混在一起，使目标组合无法审计。正确修复是恢复时间边界和状态账本，不是借未来执行信息重写信号日决策。

\section{本章产出}

\begin{outputbox}
本章应提交：冻结候选规则与并列规则；基础等权；个股、行业、流动性和换手约束配置；缓冲方案状态；约束优先级和确定性降级记录；证券目标权重、期望交易差额与现金权重；组合约束检查；自动测试与人工核对结果；以及进入第 10 章的唯一冻结 \texttt{portfolio\_version}。
\end{outputbox}

\begingroup
\small
\begin{longtable}{@{}p{4.4cm}p{4.8cm}p{5.6cm}@{}}
\caption{目标组合证券级输出 schema}\label{tab:ch09-output-schema}\\
\toprule
字段 & 含义 & 审计要求 \\
\midrule
\endfirsthead
\toprule
字段 & 含义 & 审计要求 \\
\midrule
\endhead
\texttt{security\_id}、\texttt{signal\_date} & 稳定主键与信号日 & 同一组合版本唯一 \\
\texttt{combination\_version} & 第 8 章冻结综合分版本 & 不得重选历史版本 \\
\texttt{portfolio\_version} & 目标组合版本 & 规则或约束变化即新版本 \\
\texttt{combined\_score}、\texttt{score\_rank} & 冻结分数与确定性排名 & 权重不能覆盖分数 \\
\texttt{selected}、\texttt{selection\_reason} & 入选状态与原因 & 未入选记录仍保留 \\
\texttt{base\_weight} & 约束前等权起点 & 与候选数量一致 \\
\texttt{previous\_actual\_weight} & 上一期实际证券权重 & 来自信号日已知快照 \\
\texttt{target\_weight} & 约束后目标权重 & 非负且满足硬约束 \\
\texttt{desired\_trade\_weight} & 目标减上期实际权重 & 不代表成交 \\
\texttt{industry\_asof} & 点时行业 & 生效日不晚于信号日 \\
\texttt{benchmark\_industry\_weight\_asof} & 点时基准行业权重 & 无可靠版本时使用登记替代 \\
\texttt{liquidity\_proxy\_asof} & 历史流动性代理 & 窗口截止信号日 \\
\texttt{individual\_cap} & 个股上限 & 来自冻结配置 \\
\texttt{industry\_constraint\_status} & 行业约束状态 & 保存偏离与调整 \\
\texttt{liquidity\_constraint\_status} & 流动性约束状态 & 不使用未来成交量 \\
\texttt{turnover\_constraint\_status} & 换手预算状态 & 保存收缩轨迹 \\
\texttt{portfolio\_status} & 成功、部分配置或不可行 & 不静默忽略失败 \\
\texttt{cash\_weight} & 组合级现金权重 & 与证券权重合计为 1 \\
\texttt{constraint\_version} & 约束配置版本 & 可重放优先级与降级 \\
\bottomrule
\end{longtable}
\endgroup

组合级汇总另保存证券目标权重总和、现金、目标换手、行业权重及偏离、最大个股权重、证券数量、流动性覆盖和所有失败/降级记录。

\section{章节自测}

\begin{quizbox}
\begin{enumerate}
  \item 综合得分、目标组合、期望订单、实际成交和实际持仓分别表示什么？
  \item 为什么本章选择固定数量作为主基线，而固定分位只作敏感性方案？
  \item 基础等权为什么不能直接称为最终目标权重？
  \item 截断个股权重后为什么不能无条件归一化到满仓？
  \item 行业与流动性数据的点时边界是什么？
  \item 目标换手与实际成交换手为何必须分列？
  \item 持仓缓冲如何降低交易，又引入了什么路径依赖？
  \item 哪些情况下现金权重为正是合法结果？
  \item 不可行时，确定性降级为何优于随机删股或静默放宽？
  \item 进入第 10 章前必须冻结哪些版本、状态和审计记录？
\end{enumerate}
\end{quizbox}

\section{与第 10 章《含成本回测与结果解释》的衔接}

本章交付的是冻结目标组合与期望交易差额，不是下一交易日实际持仓。第 10 章将按事件顺序生成订单、判断停牌和价格限制、模拟成交、计算费用、更新实际持仓并解释目标偏离；它必须保留无法成交的旧持仓，不能回写本章候选或目标权重。

交接包至少包含 \texttt{portfolio\_version}、证券目标权重、期望差额、现金、上一期实际持仓快照、约束状态、降级轨迹和点时数据版本。只有这些输入全部冻结后，才能开始含成本回测。

\chapter{含成本回测与结果解释}

\section*{待编写}

% TODO: 按 P03 后续里程碑要求编写本章。


\appendix

\chapter{术语与符号表}
\label{app:glossary}

\section{如何使用本附录}

本附录是第 \cref{ch:hypothesis,ch:data-timing,ch:universe-labels,ch:factor-construction,ch:preprocessing,ch:single-factor,ch:robustness,ch:factor-combination,ch:portfolio,ch:backtest} 的快速检索入口，不替代正文定义。术语按研究流程分组；“主要出现”指第一次系统解释该对象的章节，而非唯一出现位置。英文名称和代码名用于帮助读者连接公式、配置与输出字段。

若术语依赖真实数据、交易制度或团队配置，应回到对应章节的“需实际确认”说明，不得仅凭本表推断实际口径。

\section{核心研究术语}

\begingroup
\footnotesize
\setlength{\tabcolsep}{3pt}
\begin{longtable}{@{}p{3.2cm}p{9.0cm}p{2.8cm}@{}}
\caption{核心研究设计术语}\label{tab:appA-research-terms}\\
\toprule
中文名称 / 英文或代码名 & 准确解释 & 主要出现 \\
\midrule
\endfirsthead
\toprule
中文名称 / 英文或代码名 & 准确解释 & 主要出现 \\
\midrule
\endhead
研究假设\\Research hypothesis & 对冻结股票池、信息时点、可计算代理、预测标签、预期方向和证据规则的可检验陈述；它不是因子公式或交易规则。 & \cref{ch:hypothesis} \\
可证伪性\\Falsifiability & 研究设计允许预先说明哪些观察会反对原想法、触发返工或停止，而不是只接受有利结果。 & \cref{ch:hypothesis} \\
主分析\\Primary analysis & 在查看本轮结果前冻结、承担主要结论的唯一基线分析。 & \cref{ch:hypothesis} \\
预注册敏感性\\Preregistered sensitivity & 在结果前登记的少量邻近设定，用于检查结论是否依赖单一合理口径；不能用于事后选赢家。 & \cref{ch:hypothesis} \\
探索性分析\\Exploratory analysis & 查看结果后产生的新想法，只能形成下一轮假设和新版本，不能改写本轮确认性证据。 & \cref{ch:hypothesis} \\
样本内\\In-sample & 研究者用于形成、实现或有限选择方案的时间段；其证据不能冒充最终样本外证据。 & \cref{ch:robustness} \\
验证集\\Validation set & 在时间顺序下用于有限、预注册选择的区间；每次访问和选择理由均应留痕。 & \cref{ch:robustness} \\
测试集\\Test set & 用于评价已经冻结方案的后续时间区间；看过结果后再改方案会消耗其测试身份。 & \cref{ch:robustness} \\
最终留出集\\Final holdout & 全部研究设计锁定后原则上只查看一次的最后时间区间；重复访问必须披露。 & \cref{ch:robustness} \\
数据窥探\\Data snooping & 反复查看同一数据的收益、RankIC、覆盖或成本等结果并据此调整方案，使偶然模式被当作证据。 & \cref{ch:robustness} \\
选择性汇报\\Selective reporting & 只展示有利参数、样本或成功实验而隐藏失败、报错和证据不足，导致选择次数不可见。 & \cref{ch:robustness} \\
\bottomrule
\end{longtable}
\endgroup

\section{数据与时间术语}

\begingroup
\footnotesize
\setlength{\tabcolsep}{3pt}
\begin{longtable}{@{}p{3.2cm}p{9.0cm}p{2.8cm}@{}}
\caption{数据与时间术语}\label{tab:appA-data-time-terms}\\
\toprule
中文名称 / 英文或代码名 & 准确解释 & 主要出现 \\
\midrule
\endfirsthead
\toprule
中文名称 / 英文或代码名 & 准确解释 & 主要出现 \\
\midrule
\endhead
点时数据\\Point-in-Time, PIT & 对每个历史信号截止，只使用当时已经可得且可追溯版本的数据；“数据库今天能查到”不代表历史当时可见。 & \cref{ch:data-timing} \\
报告期\\Reporting period & 财务或业务数据所描述的经济期间，不等于信息向市场公开或研究者可用的日期。 & \cref{ch:data-timing} \\
公告日\\Announcement date & 信息对外公布的日期或时间戳；若只有日期而无时刻，是否赶上当日信号截止仍需确认。 & \cref{ch:data-timing} \\
可得日\\Availability date & 数据经过发布、采集和处理后能够进入研究系统的最早时点，可能晚于公告日。 & \cref{ch:data-timing} \\
生效日\\Effective date & 成分、行业、证券状态或规则开始适用于研究对象的日期；连接时应检查有效区间。 & \cref{ch:data-timing} \\
信号截止\\Signal cutoff & 本期因子和目标组合允许读取信息的最晚时点；截止后的状态不能回写信号横截面。 & \cref{ch:data-timing} \\
信号日\\Signal date & 冻结股票池、因子与目标组合的日期，本案例为每周最后一个交易日收盘后。 & \cref{ch:universe-labels} \\
成交日\\Execution date & 信号日之后按配置尝试模拟成交的交易日；其状态只在执行阶段读取。 & \cref{ch:universe-labels} \\
持有期\\Holding period & 从标签或实际持仓的起始成交代理时点到下一结束时点的区间；端点和日历必须版本化。 & \cref{ch:universe-labels} \\
标签窗口\\Label window & 用于计算未来收益标签的起止时点；起点必须严格晚于信号截止。 & \cref{ch:universe-labels} \\
向后点时连接\\as-of join & 对每个左表时点选择不晚于它的最近合法右表记录，并同时核验实体键、可得时间和有效区间。 & \cref{ch:data-timing} \\
前视偏差\\Look-ahead bias & 特征、样本选择或执行假设使用了历史当时尚不可得的信息，使回测获得未来知识。 & \cref{ch:data-timing} \\
幸存者偏差\\Survivorship bias & 用今天仍存在或仍在指数中的证券重建历史，遗漏后来退市、调出或更名的对象。 & \cref{ch:universe-labels} \\
历史修订版本\\Historical revision version & 同一观察在不同时间发布的版本；研究应保存当时版本或明确无法恢复历史修订的限制。 & \cref{ch:data-timing} \\
稳定证券标识\\Stable security identifier & 跨代码变化保持同一实体、又能避免代码复用误连的证券主键；展示代码不能自动替代它。 & \cref{ch:universe-labels} \\
\bottomrule
\end{longtable}
\endgroup

\section{因子与评价术语}

\begingroup
\footnotesize
\setlength{\tabcolsep}{3pt}
\begin{longtable}{@{}p{3.2cm}p{9.0cm}p{2.8cm}@{}}
\caption{因子构造、预处理与评价术语}\label{tab:appA-factor-evaluation-terms}\\
\toprule
中文名称 / 英文或代码名 & 准确解释 & 主要出现 \\
\midrule
\endfirsthead
\toprule
中文名称 / 英文或代码名 & 准确解释 & 主要出现 \\
\midrule
\endhead
原始因子\\Raw factor, \texttt{raw\_value} & 依据冻结公式和点时输入直接计算的因子值，尚未经过第 5 章截面预处理。 & \cref{ch:factor-construction} \\
处理后因子\\\texttt{processed\_value} & 在同一信号日截面完成有效筛选、缩尾、方向统一、标准化及登记暴露处理后的数值。 & \cref{ch:preprocessing} \\
有效值\\Valid value & 满足因子定义、输入可得、分母合法且数值有限的观察；有效但极端不等于错误。 & \cref{ch:factor-construction} \\
缺失值\\Missing value & 输入或结果没有可用数值的状态；缺失记录应保留原因，不能无条件填零。 & \cref{ch:preprocessing} \\
风险原因\\Risk reason / risk flag & 对仍可计算但需要解释的会计、价格、暴露或交易风险做标记；它与导致因子无效的 \texttt{invalid\_reason} 不同。 & \cref{ch:factor-construction} \\
缩尾\\Winsorization & 将超出冻结截面边界的有效极端值压到边界，不删除证券，也不使用全样本未来分布。 & \cref{ch:preprocessing} \\
绝对中位差\\Median absolute deviation, MAD & 当期截面中各值与中位数绝对距离的中位数，用于构造稳健的缩尾尺度。 & \cref{ch:preprocessing} \\
标准分\\z-score & 当期有效截面内减去均值并除以标准差得到的无量纲数值；退化截面应返回状态。 & \cref{ch:preprocessing} \\
截面排名\\Cross-sectional rank & 在同一信号日和同一因子内按冻结并列规则排序，主要保留相对顺序。 & \cref{ch:preprocessing} \\
行业内处理\\Within-industry processing & 在点时行业分组内部排名、去均值或标准化，使证券更多地与同业比较。 & \cref{ch:preprocessing} \\
残差化\\Residualization & 用当期行业、市值等解释因子截面差异，并把回归残差作为替代处理版本；设计矩阵退化必须留痕。 & \cref{ch:preprocessing} \\
中性化\\Neutralization & 降低指定非目标暴露的统称，可由行业内处理、残差化或组合约束实现；不是所有因子的机械必选项。 & \cref{ch:preprocessing} \\
因子方向\\Factor direction & 冻结因子数值与预期未来相对收益的方向符号；评价阶段不得按结果临时翻向。 & \cref{ch:factor-construction} \\
因子版本\\\texttt{factor\_version} & 标识公式、输入语义和实现身份；实质变化产生新版本，不覆盖历史输出。 & \cref{ch:factor-construction} \\
覆盖率\\Coverage & 有效因子证券数占冻结信号池证券数的比例，分母不能改成“已有值样本”。 & \cref{ch:preprocessing} \\
可评价覆盖率\\Evaluable coverage & 同时具有冻结因子值和冻结标签的证券数占信号池的比例，用于披露连接标签后的样本损失。 & \cref{ch:single-factor} \\
信息系数\\Information coefficient, IC & 同一信号日可评价横截面中因子值与未来收益标签的 Pearson 相关。 & \cref{ch:single-factor} \\
秩信息系数\\Rank information coefficient, RankIC & 同一信号日因子排名与未来收益排名的相关，评价排序方向而非可交易收益。 & \cref{ch:single-factor} \\
分组收益\\Grouped return & 按当期因子排序和冻结分组规则汇总各组未来标签；诊断头尾差不等于真实策略收益。 & \cref{ch:single-factor} \\
单调性\\Monotonicity & 因子组从低到高时，未来收益是否大体按预期方向变化；不能只检查头尾两组。 & \cref{ch:single-factor} \\
因子自相关\\Factor autocorrelation & 同一证券因子值或排名在相邻信号日的相关，用于描述信号持续性。 & \cref{ch:single-factor} \\
因子换手\\Factor turnover & 因子排名或诊断分组成员变化的代理，不等同于目标、订单或实际成交换手。 & \cref{ch:single-factor} \\
暴露\\Exposure & 因子、综合分或组合对行业、市值、流动性、波动等非目标维度的系统关系；暴露需要解释但不自动等于错误。 & \cref{ch:single-factor} \\
稳健性\\Robustness & 证据对时间推进、合理参数邻域、股票池、子样本和预注册口径变化的抵抗力。 & \cref{ch:robustness} \\
参数尖峰\\Parameter spike & 只有一个参数点表现突出、相邻合理设定不支持同样结论的形态，提示选择偏差或不稳定。 & \cref{ch:robustness} \\
条件分组\\Conditional sorting & 先按一个冻结因子分组，再在组内评价另一个因子，以检查其条件增量信息。 & \cref{ch:factor-combination} \\
留一消融\\Leave-one-out ablation & 在严格共同样本和同一评价规则下依次移除一个因子，比较完整组合与留一版本。 & \cref{ch:factor-combination} \\
增量价值\\Incremental value & 新因子在已有因子和共同暴露之外，是否仍提供可解释、稳健且基本可交易的信息。 & \cref{ch:factor-combination} \\
\bottomrule
\end{longtable}
\endgroup

\section{组合与执行术语}

\begingroup
\footnotesize
\setlength{\tabcolsep}{3pt}
\begin{longtable}{@{}p{3.2cm}p{9.0cm}p{2.8cm}@{}}
\caption{组合、成交与回测术语}\label{tab:appA-portfolio-backtest-terms}\\
\toprule
中文名称 / 英文或代码名 & 准确解释 & 主要出现 \\
\midrule
\endfirsthead
\toprule
中文名称 / 英文或代码名 & 准确解释 & 主要出现 \\
\midrule
\endhead
综合得分\\Combined score, \(S_{i,t}\) & 对方向和尺度统一、正式纳入且共同有效的因子按冻结规则合成的证券层分数；不是目标权重。 & \cref{ch:factor-combination} \\
目标权重\\\texttt{target\_weight} & 在信号日点时信息和组合约束下希望持有的证券权重；成交失败前仍只是目标。 & \cref{ch:portfolio} \\
上一期实际权重\\\path{previous_actual_weight} & 上一调仓完成后延续到当前信号日、已知的实际证券权重快照。 & \cref{ch:portfolio} \\
期望交易差额\\\path{desired_trade_weight} & 目标权重减上一期实际权重，正值表示期望增加、负值表示期望减少；不是订单或成交。 & \cref{ch:portfolio} \\
订单\\Order & 由期望差额、冻结净值、交易单位和取整规则形成的买卖请求；订单存在不代表已经成交。 & \cref{ch:backtest} \\
成交\\Fill / execution & 在成交日状态和价格代理下真正执行的数量与价格事件，并关联原订单。 & \cref{ch:backtest} \\
部分成交\\Partial fill & 请求量只有一部分执行；成交量、未成交量与请求量必须守恒。 & \cref{ch:backtest} \\
未成交\\Unfilled & 请求量没有执行；买入失败保留现金，卖出失败保留旧持仓。 & \cref{ch:backtest} \\
实际持仓\\Actual holding & 由旧持仓、实际成交和公司行动按事件顺序共同决定的证券与现金状态。 & \cref{ch:backtest} \\
现金权重\\Cash weight & 现金占组合净值的比例；目标阶段的现金与实际成交后现金应分开记录。 & \cref{ch:portfolio} \\
目标换手\\Target turnover & 目标权重相对上一期实际权重变化的单边规模，仍未考虑订单取整或成交失败。 & \cref{ch:portfolio} \\
订单换手\\Order turnover & 实际送入成交模拟的订单规模占组合净值的比例，可能因取整和现金规则偏离目标换手。 & \cref{ch:backtest} \\
实际成交换手\\Realized turnover & 仅以实际成交金额计算的换手，未成交订单不进入分子。 & \cref{ch:backtest} \\
毛收益\\Gross return & 实际持仓在扣除本期交易成本前的组合收益，需与公司行动和现金口径一致。 & \cref{ch:backtest} \\
净收益\\Net return & 毛收益减去本期成本占净值比例后的结果；费用不能在现金和收益分解中重复扣除。 & \cref{ch:backtest} \\
交易成本\\Transaction cost & 显性费用、滑点和冲击等因实际成交产生的成本集合；历史口径需实际确认。 & \cref{ch:backtest} \\
滑点\\Slippage & 模拟执行价相对冻结基准价格的不利偏离代理，方向和价格基准需配置。 & \cref{ch:backtest} \\
冲击成本\\Market impact & 订单相对流动性规模可能引起的额外价格代价；本课程仅使用透明配置化概念模型。 & \cref{ch:backtest} \\
最大回撤\\Maximum drawdown & 净值相对此前历史峰值的最大下跌幅度，必须基于完整且口径一致的序列。 & \cref{ch:backtest} \\
夏普比率\\Sharpe ratio & 超过冻结无风险收益口径的平均收益相对波动的比率；不能作为单一研究结论。 & \cref{ch:backtest} \\
跟踪误差\\Tracking error & 组合相对冻结基准的超额收益波动，频率、年化与缺失期处理需一致。 & \cref{ch:backtest} \\
信息比率\\Information ratio, IR & 平均超额收益相对跟踪误差的比率，依赖同一基准与样本区间。 & \cref{ch:backtest} \\
目标—实际偏差\\Target--actual gap & 实际权重减目标权重，用于解释成交失败、旧持仓延续、现金和约束导致的实现差异。 & \cref{ch:backtest} \\
\bottomrule
\end{longtable}
\endgroup

\section{数学符号表}

\begingroup
\footnotesize
\setlength{\tabcolsep}{3pt}
\begin{longtable}{@{}p{3.2cm}p{9.0cm}p{2.8cm}@{}}
\caption{主体讲义主要数学符号}\label{tab:appA-symbols}\\
\toprule
符号 / 代码名 & 含义 & 主要出现 \\
\midrule
\endfirsthead
\toprule
符号 / 代码名 & 含义 & 主要出现 \\
\midrule
\endhead
\(t\) & 周度信号日或对应的信号截面索引。 & \cref{ch:universe-labels} \\
\(\tau(t)\) & 与信号日 \(t\) 对应的下一交易日成交代理时点。 & \cref{ch:universe-labels} \\
\(\mathcal{U}^{\mathrm{research}}_t\) & 信号日按历史生效成分构建的研究股票池。 & \cref{ch:universe-labels} \\
\(\mathcal{U}^{\mathrm{signal}}_t\) & 仅用信号截止及以前信息过滤后的信号有效股票池。 & \cref{ch:universe-labels} \\
\(\mathcal{V}_{f,t}\) & 因子 \(f\) 在 \(t\) 日满足定义且数值有限的有效证券集合。 & \cref{ch:preprocessing} \\
\(\mathcal{V}^{\mathrm{joint}}_t\) & 正式纳入因子在 \(t\) 日的共同有效证券交集。 & \cref{ch:factor-combination} \\
\(\mathcal{C}_t\) & 经综合分、排序和组合规则形成的目标组合候选集合。 & \cref{ch:universe-labels} \\
\(\mathcal{X}_{\tau(t)}\) & 成交时点依据当时状态和价格代理判断的可执行证券集合。 & \cref{ch:universe-labels} \\
\(\mathcal{H}_{\tau(t)}\) & 订单处理后真正持有的证券集合，包括未能卖出的旧持仓。 & \cref{ch:universe-labels} \\
\(f_{i,t}\) & 证券 \(i\) 在信号日 \(t\) 的冻结因子值。 & \cref{ch:single-factor} \\
\(y_{i,t}\) & 证券 \(i\) 对应冻结标签窗口的未来收益。 & \cref{ch:universe-labels} \\
\(z_{f,i,t}\) & 因子 \(f\) 对证券 \(i\) 在 \(t\) 日方向和尺度统一后的处理值。 & \cref{ch:factor-combination} \\
\(S_{i,t}\) & 证券 \(i\) 在 \(t\) 日的冻结多因子综合得分。 & \cref{ch:factor-combination} \\
\(w^{\mathrm{base}}_{i,t}\) & 约束调整前、由候选规则产生的基础目标权重。 & \cref{ch:portfolio} \\
\(w^{\mathrm{target}}_{i,t}\) & 完成信号日约束与降级规则后的证券目标权重。 & \cref{ch:portfolio} \\
\(h^{\mathrm{actual}}_{i,t}\) & 证券 \(i\) 的实际持仓权重；代码和第 10 章公式也使用 \texttt{actual\_weight} 或 \(w^{\mathrm{actual}}_{i,t}\)。 & \cref{ch:backtest} \\
\(w^{\mathrm{cash}}_t\) & 目标或实际现金占组合净值的权重，使用时须注明层次。 & \cref{ch:portfolio} \\
\(\mathrm{IC}_t\) & 信号日 \(t\) 因子值与未来收益的横截面 Pearson 相关。 & \cref{ch:single-factor} \\
\(\mathrm{RankIC}_t\) & 信号日 \(t\) 因子排名与未来收益排名的横截面相关。 & \cref{ch:single-factor} \\
\(\mathrm{Coverage}_{f,t}\) & 因子 \(f\) 在 \(t\) 日有效证券数除以冻结信号池证券数。 & \cref{ch:preprocessing} \\
\(\mathrm{Turnover}_t\) & 某一明确层次在 \(t\) 的换手；必须用上标或字段名区分因子、目标、订单和成交。 & \cref{ch:single-factor,ch:portfolio,ch:backtest} \\
\(\mathrm{GrossReturn}_t\) & 基于实际持仓、扣除本期成本前的组合收益。 & \cref{ch:backtest} \\
\(\mathrm{NetReturn}_t\) & 毛收益扣除本期成本占净值比例后的组合收益。 & \cref{ch:backtest} \\
\bottomrule
\end{longtable}
\endgroup

\section{状态码与版本字段}

\begingroup
\footnotesize
\setlength{\tabcolsep}{3pt}
\begin{longtable}{@{}p{3.5cm}p{8.7cm}p{2.8cm}@{}}
\caption{常用状态码与版本字段}\label{tab:appA-status-version}\\
\toprule
字段或状态 & 用途 & 主要出现 \\
\midrule
\endfirsthead
\toprule
字段或状态 & 用途 & 主要出现 \\
\midrule
\endhead
\texttt{is\_valid} / \texttt{invalid\_reason} & 区分因子是否满足定义，并保留历史不足、输入缺失、非法分母等原因；无效记录不能静默删除。 & \cref{ch:factor-construction} \\
\texttt{processing\_status} & 记录预处理成功、输入无效、MAD 或标准差退化、残差化失败等状态。 & \cref{ch:preprocessing} \\
\texttt{OK}, \texttt{INSUFFICIENT\_N}, \texttt{CONSTANT\_FACTOR}, \texttt{CONSTANT\_LABEL}, \texttt{MANY\_TIES} & 单信号日评价的典型状态，使不能计算 IC 或分组的原因可追溯；具体阈值来自配置。 & \cref{ch:single-factor} \\
\texttt{PASS / FAIL /}\\\texttt{INSUFFICIENT\_EVIDENCE} & 模板中的空白判断枚举，分别表示通过、失败和证据不足；必须由实际证据填写。 & \cref{ch:hypothesis} \\
\texttt{PARTIAL\_ALLOCATION} / \texttt{INFEASIBLE} & 目标组合部分配置或约束不可行状态；前者可保留现金，后者不得静默放宽硬约束。 & \cref{ch:portfolio} \\
\texttt{order\_status} / \texttt{order\_reason} & 说明订单是否有效、被取整或受现金等规则影响；订单不是成交。 & \cref{ch:backtest} \\
\texttt{fill\_status} / \texttt{unfilled\_reason} & 记录完全、部分、未成交、未知或无效及原因；不得回写信号日对象。 & \cref{ch:backtest} \\
\texttt{data\_version} / \texttt{universe\_version} / \texttt{label\_version} & 标识数据快照、历史股票池和标签规则；不得只写“最新”。 & \cref{ch:data-timing,ch:robustness} \\
\texttt{factor\_version} / \texttt{preprocess\_version} & 标识因子公式和预处理流程；评价与合成只能读取已冻结版本。 & \cref{ch:factor-construction,ch:preprocessing} \\
\texttt{combination\_version} & 标识正式因子集合、共同样本与综合分规则；因子集合或权重变化即新版本。 & \cref{ch:factor-combination} \\
\texttt{portfolio\_version} / \texttt{constraint\_version} & 标识目标组合规则、约束优先级和降级逻辑。 & \cref{ch:portfolio} \\
\texttt{execution\_version} / \texttt{holding\_version} / \texttt{backtest\_version} & 标识成交规则、持仓账本和完整回测配置，使订单、成交、现金和收益可重放。 & \cref{ch:backtest} \\
\bottomrule
\end{longtable}
\endgroup

\section{容易混淆概念对照}

\begingroup
\footnotesize
\setlength{\tabcolsep}{3pt}
\begin{longtable}{@{}p{4.0cm}p{5.4cm}p{5.6cm}@{}}
\caption{容易混淆的对象及其职责边界}\label{tab:appA-confusable-objects}\\
\toprule
对象组 & 各自回答的问题 & 不得混用的原因 \\
\midrule
\endfirsthead
\toprule
对象组 & 各自回答的问题 & 不得混用的原因 \\
\midrule
\endhead
报告期 / 公告日 / 可得日 & 报告期描述经济期间；公告日描述对外发布；可得日描述研究系统最早可用。 & 按报告期直接前填会让尚未发布的信息进入历史信号。 \\
原始因子 / 处理后因子 / 综合分 & 原始因子是冻结公式输出；处理后因子完成当日截面处理；综合分合成多个正式因子。 & 三者版本、尺度和用途不同，综合分也不是目标持仓。 \\
研究池 / 信号池 / 目标候选 / 实际持仓 & 研究池回答历史成分；信号池回答谁可比较；候选回答想持有谁；实际持仓回答执行后真正持有什么。 & 成交日状态不能回删信号池，未能卖出的证券也可能不在目标候选但仍在实际持仓。 \\
目标权重 / 期望交易 / 订单 / 成交 / 实际持仓 & 依次表示想持多少、相对旧持仓希望改变多少、发出什么请求、实际执行多少、最终留下什么。 & 把目标当成交会隐藏部分成交、未成交、现金和旧持仓延续。 \\
因子换手 / 目标换手 / 订单换手 / 实际成交换手 & 分别描述排名或成员变化、目标权重变化、订单规模和实际成交规模。 & 分子、分母与发生阶段不同，不能共用一个“turnover”列。 \\
毛收益 / 成本 / 净收益 & 毛收益是实际持仓扣成本前结果；成本来自实际成交；净收益为毛收益扣成本。 & 成本若已进入现金账本，不得在收益分解中重复扣除。 \\
数据问题 / 不利结果 / 证据不足 & 数据问题是可定位的时序、版本或实现错误；不利结果是可信证据不支持假设；证据不足是样本或覆盖不足以评价。 & 只有数据问题允许返工；不利结果不能借返工继续搜索，证据不足也不能记为通过。 \\
\bottomrule
\end{longtable}
\endgroup

\chapter{研究检查清单}

\section*{待编写}

% TODO: 后续里程碑补充检查清单。

\chapter{研究模板}

\section*{待编写}

% TODO: 后续里程碑补充研究模板。

\include{appendices/appD_production_overview}
\chapter{AI 使用说明}
\label{app:ai-usage}

\section*{如何使用本附录}

本附录依据仓库中实际存在的 \path{ai_usage/raw_prompts.md} 与 \texttt{ai\_usage/AI使用记录.md} 整理，只提供角色边界、核验状态和索引。完整长 prompt 不复制进主 PDF，而与主讲义并列提交。若本附录与原始日志不一致，应以版本化原始日志为准并记录修订。

\section{AI 在本项目中的角色}

AI 在本项目中承担结构、草稿和审查辅助：帮助拆解任务、收敛课程范围、组织教学顺序、生成 LaTeX 初稿、执行静态扫描、修复已定位的编译问题，并整理错误与一致性清单。AI 输出始终是待核验材料，不拥有研究决策权。

AI 不是本项目的真实数据提供者、实证结果提供者、投资建议提供者、最终技术责任人或人工复核替代者。日志中的“提交”“静态检查”或“代码检查”也不自动等于内容已被作者采用。

\section{哪些工作由 AI 辅助}

现有日志显示，AI 辅助工作主要包括：

\begin{itemize}
  \item 任务拆解与范围收敛：把广泛蓝图映射到冻结的十章主体；
  \item 教学结构设计：为点时数据、股票池、因子、评价、组合与回测建立前向接口；
  \item LaTeX 初稿：撰写章节、图表、代码清单、模板和附录；
  \item 静态审查：检查章节顺序、字段覆盖、环境配对、标签/引用、禁止命令和构建产物；
  \item 编译修复辅助：根据用户提供的 Overleaf 错误定位 tcolorbox、TikZ 和数学转写问题；
  \item 错误清单：把前视、幸存者偏差、样本选择、账本不守恒等风险转为测试与停止线；
  \item 可读性与一致性检查：统一术语、版本字段、交叉引用和阶段边界。
\end{itemize}

这些工作没有产生真实金融数据，也没有验证三个教学因子的实证有效性。

\section{哪些决定必须由人工负责}

作者或指定复核人必须负责：课程范围与最终采用；数据来源、授权和字段语义；历史指数与证券主键；财务、价格、复权、公司行动和成交口径；研究阈值与停止条件；代码与公式正确性；真实数据运行；回测解释；是否满足交付要求；以及任何对外使用或投资相关决定。

人工还必须审查 AI 是否遗漏上下文、误读接口、引入未经定义对象、制造过密版面或把示意写成事实。AI 不能自我签署“人工核验完成”，也不能用模型内部推理替代可公开的验证证据。

\section{提问如何逐步收敛}

现有记录从 P01 的首轮培训蓝图开始，P02 将广泛课程收敛为十章主体、贯穿案例和标准交付结构，P03 再把用户提供的正式蓝图纳入仓库。随后按研究事件顺序推进：P03-1 建立数据时间边界，P04—P06 完成股票池、因子和预处理，P07—P08 处理单因子章节及 LaTeX 阻塞，P08-ROBUSTNESS 至 P11 完成稳健性、合成、组合和含成本回测，P12 回补第 1 章研究假设，P13 补齐附录与提交说明，P14 更新结项 README，P15 从用户提供的历史对话补录 P01 与 P02，P16 提交并推送结项记录修改。

每轮 prompt 越来越明确地冻结允许修改文件、章节职责、禁止事项、验收和 Git 操作。这种收敛减少了章节间职责漂移，但也意味着早期蓝图的广泛生产化范围不能直接覆盖后来的十章冻结结构；冲突均按当前任务和 \path{AGENTS.md} 的优先级处理。

\section{事实、公式、代码和编译核验}

核验分为不同层次，不能相互替代：

\begin{itemize}
  \item \textbf{事实与制度：}应由官方规则、真实数据字典或责任人确认；日志没有证明任何实际指数、费用或交易制度口径。
  \item \textbf{公式与接口：}通过正文交叉核对、匿名手算、边界条件和代码职责检查；仍需作者复核数学与业务语义。
  \item \textbf{代码：}现有记录主要检查静态结构、禁止字段、确定性、AST 或伪代码接口；仓库当前没有真实数据研究 Notebook 或完整研究代码可供端到端验证。
  \item \textbf{LaTeX：}多条记录明确说明本地没有 XeLaTeX/latexmk；虽然用户任务描述当前 PDF 可编译，仓库日志没有可据此改写为“本地已编译通过”的证据。
  \item \textbf{PDF 视觉：}若能生成 PDF，应渲染逐页检查表格、TikZ、中文、页眉页脚与分页；当前日志未记录完成该人工视觉验收。
\end{itemize}

\section{已知限制与未完成核验}

当前 AI 日志存在以下明确状态：

\begin{itemize}
  \item P01 与 P02 已根据用户提供的历史 ChatGPT 对话补录原始 prompt、附件名称和页面可见输出摘要；历史输出文件未在当前仓库单独保存，且尚未与当前蓝图逐字比对；
  \item P08 包含一个编译修复记录和一个稳健性章节记录，两个编号均保留，不合并或重命名；
  \item 既有原始 prompt 按任务追加，未检测到完全相同的重复任务块；不同任务重复强调 Git、安全和编译要求属于真实原文，不删除；
  \item 既有 AI 使用记录均为“待人工确认”或“人工核验尚未完成”，没有证据支持改写为已审核；
  \item 多数记录明确说明本地缺少 TeX，实际 PDF 页数、版面和交叉引用仍需统一审查；
  \item 仓库当前没有真实金融数据、研究 Notebook、源代码实现、配置文件或实证输出，主体公式和伪代码不能被描述为已验证策略。
\end{itemize}

上述限制只做披露。补录记录必须保留来源和未核验边界，不把页面未显示的模型、绝对日期、文件全文、人工签署、数据结果或编译记录补造成事实。

\section{AI 提问索引}

\begingroup
\scriptsize
\setlength{\tabcolsep}{1pt}
\begin{longtable}{@{}p{1.35cm}p{2.05cm}p{2.15cm}p{2.35cm}p{2.25cm}p{1.75cm}p{1.45cm}p{1.85cm}@{}}
\caption{基于仓库真实日志的 AI 使用索引}\label{tab:appE-ai-index}\\
\toprule
\texttt{prompt\_id} & 任务名称 & 提问目的 & 主要输出 & 修改文件 & 验证方式 & 人工核验状态 & 已知限制 \\
\midrule
\endfirsthead
\toprule
\texttt{prompt\_id} & 任务名称 & 提问目的 & 主要输出 & 修改文件 & 验证方式 & 人工核验状态 & 已知限制 \\
\midrule
\endhead
\path{P01-BLUEPRINT-001} & 首轮培训蓝图 & 定义能力、流程、边界和课程框架 & 19 章首轮蓝图与验收框架 & 历史输出文件 & 对话页面核对 & 文件全文待核验 & 原文件未入库 \\
\path{P02-SCOPE-001} & 范围收敛与重构 & 收敛主体并设计贯穿案例 & 十章结构、三因子案例与质量闸门 & 历史输出文件、后续蓝图 & 对话与结构核对 & 文件全文待核验 & 附件原件未入库 \\
\path{P03-BLUEPRINT-001} & 正式蓝图入库 & 保存用户蓝图并建立索引 & 蓝图与文档索引 & \path{docs/}、AI 日志 & 文件与 Git 检查 & 待人工确认 & 蓝图范围较广 \\
\path{P03-1-001} & 第 2 章 & 建立数据与时间边界 & PIT、时序、字段审计与代码示例 & 第 2 章、AI 日志 & 静态检查 & 尚未完成 & 无本地 TeX、无真实数据 \\
\path{P04-001} & 第 3 章 & 冻结股票池与标签边界 & 五类集合、标签与边界测试 & 第 3 章、AI 日志 & 静态检查 & 尚未完成 & 制度与价格代理待确认 \\
\path{P05-001} & 第 4 章 & 实现三类原始因子 & 定义卡、公式、schema 与核对 & 第 4 章、AI 日志 & 静态与人工样例计划 & 尚未完成 & 财务与价格口径待确认 \\
\path{P06-001} & 第 5 章 & 建立截面预处理基线 & MAD、方向、标准化与暴露 & 第 5 章、AI 日志 & 静态检查 & 尚未完成 & 参数与行业口径待确认 \\
\path{P07-001} & 修复与第 6 章 & 修复 TikZ/页眉并建立单因子评价 & 工程修复、五维证据链 & 工程文件、第 2—6 章、AI 日志 & 静态扫描与 Git & 尚未完成 & 无本地 TeX、无真实数据 \\
\path{P08-001} & LaTeX 阻塞修复 & 修复裸标签和数学转写 & tcolorbox、TikZ 与公式修复 & 第 3—6 章、AI 日志 & 环境、标签和引用扫描 & 尚未完成 & PDF 未在本地验证 \\
\path{P08-ROBUSTNESS-001} & 第 7 章 & 建立样本外与反过拟合流程 & 时间切分、参数邻域与实验台账 & 第 7 章、AI 日志 & 静态与结构检查 & 尚未完成 & 样本、阈值待确认 \\
\path{P09-001} & 第 8 章 & 检查冗余并合成因子 & 共同覆盖、条件分组、消融与等权分 & 第 8 章、AI 日志 & 静态、算术与接口检查 & 尚未完成 & 无真实增量结果 \\
\path{P10-001} & 第 9 章 & 从综合分生成约束目标 & 目标权重、约束、现金与降级 & 第 9 章、AI 日志 & 静态、AST 与约束检查 & 尚未完成 & 资金与约束待确认 \\
\path{P11-001} & 第 10 章 & 建立含成本事件回测 & 订单、成交、持仓、现金和收益账本 & 第 10 章、进度、AI 日志 & 静态、账本与接口检查 & 尚未完成 & 执行、费用、基准待确认 \\
\path{P12-001} & 第 1 章 & 把想法转为可证伪任务 & 假设卡、预注册、停止线与研究地图 & 第 1 章、进度、AI 日志 & 结构、引用与 AST 检查 & 尚未完成 & 无本地 TeX、阈值待确认 \\
\path{P13-A} & 附录 A—C & 建立检索、清单和空白模板 & 术语表、99 项清单、14 份模板 & 附录 A—C、AI 日志 & 字段、标签、引用和结构检查 & 尚未完成 & PDF 页数和版面待验 \\
\path{P13-B} & 附录 D—E & 建立生产概览、AI 说明和提交索引 & 生产流程、日志索引与提交说明 & 附录 D—E、进度、提交说明、AI 日志 & 日志比对与静态检查 & 尚未完成 & 无本地 TeX、人工复核待完成 \\
\path{P14-README-001} & 结项 README & 更新项目状态与验收入口 & 状态、编译、边界与结项清单 & README、AI 日志 & 链接与静态检查 & 尚未完成 & PDF 与代码未交付 \\
\path{P15-RESTORE-001} & 补录 P01/P02 & 恢复历史提问与来源 & 两条原始提问、索引和状态修正 & AI 日志、附录 E、提交说明 & 历史页面与静态核对 & 文件全文待核验 & 历史原件未入库 \\
\path{P16-GIT-001} & 提交结项记录 & 提交并推送当前修改 & README、补录与索引提交 & README、AI 日志、附录 E & Git 与静态检查 & 内容验收待完成 & 无本地 TeX \\
\bottomrule
\end{longtable}
\endgroup

表中“静态检查”仅表示检查了文本或代码结构，不表示实证、编译或人工核验通过。P01 与 P02 是根据用户提供的历史对话补录，来源与未核验项记录在两份 Markdown 日志中。P13-B 与 P13-A 来自同一条完整用户任务；完整原文只在 \path{raw_prompts.md} 保存一次，避免在日志和主 PDF 中重复膨胀。

\section{完整原始提问如何提交}

完整原始提问按任务顺序保存在 \path{ai_usage/raw_prompts.md}，AI 产出摘要、修改文件、验证方式、人工核验状态和限制保存在 \texttt{ai\_usage/AI使用记录.md}。两者均应与主 PDF 和源码仓库并列提交，并保持版本一致。

主 PDF 只包含本索引，不复制数百行 prompt。若需要审查具体任务，应使用 \texttt{prompt\_id} 在两份 Markdown 日志中定位；发现记录缺失时应标记“记录缺失”或“待人工确认”，不能凭记忆重写。

\section{最终责任声明}

AI 参与了范围整理、结构设计、LaTeX 草稿、静态审查和错误修复辅助。作者负责范围决定、内容采用、事实与公式核验、代码和数据验证、人工复核、编译与最终提交。

本教学材料不构成投资建议。匿名案例、空白模板和结构示意不代表真实证券、真实数据库或实证结果；未经真实数据运行和独立核验的部分不得被描述为已验证策略。


\end{document}
